\documentclass[11pt,a4paper]{article}
\usepackage[T1]{fontenc}
\usepackage[utf8]{inputenc}
\usepackage{lmodern,microtype}
\usepackage[a4paper,margin=25mm,headheight=15pt]{geometry}
\usepackage{amsmath,amssymb,amsthm,mathtools}
\usepackage{booktabs,array,longtable,enumitem}
\usepackage[dvipsnames]{xcolor}
\usepackage{fancyhdr,titlesec}
\usepackage{xurl}
\usepackage[unicode,colorlinks=true,linkcolor=MidnightBlue,citecolor=MidnightBlue,urlcolor=MidnightBlue]{hyperref}
\usepackage{bookmark}
\definecolor{ink}{RGB}{27,49,70}
\titleformat{\section}{\large\bfseries\color{ink}}{\thesection}{0.7em}{}
\titleformat{\subsection}{\normalsize\bfseries\color{ink}}{\thesubsection}{0.7em}{}
\pagestyle{fancy}\fancyhf{}
\fancyhead[L]{\small Finite-core feedback transseries}
\fancyhead[R]{\small Research draft}
\fancyfoot[C]{\thepage}
\renewcommand{\headrulewidth}{0.3pt}
\setlength{\parindent}{1.2em}\setlength{\parskip}{3pt}
\setlength{\emergencystretch}{2em}
\setlength{\belowcaptionskip}{6pt}
\setlist{itemsep=3pt,topsep=5pt,leftmargin=1.7em}
\numberwithin{equation}{section}
\newtheorem{theorem}{Theorem}[section]
\newtheorem{proposition}[theorem]{Proposition}
\newtheorem{lemma}[theorem]{Lemma}
\newtheorem{corollary}[theorem]{Corollary}
\theoremstyle{definition}
\newtheorem{definition}[theorem]{Definition}
\newtheorem{example}[theorem]{Example}
\newtheorem{conjecture}[theorem]{Conjecture}
\theoremstyle{remark}
\newtheorem{remark}[theorem]{Remark}
% ed. (2026-09-29): unnumbered environment for ProveIt editorial notes, so that the article's own numbering is unchanged
\newtheorem*{ednote}{Editorial note (ProveIt, 2026-09-29)}
\newcommand{\N}{\mathbb N}\newcommand{\R}{\mathbb R}\newcommand{\C}{\mathbb C}
\newcommand{\ee}{\mathrm e}\newcommand{\dd}{\mathrm d}
\newcommand{\Prob}{\mathbb P}\newcommand{\E}{\mathbb E}
\newcommand{\Normal}{\mathcal N}\newcommand{\Acal}{\mathcal A}
\newcommand{\Uhat}{\widehat U}\newcommand{\Vhat}{\widehat V}
\newcommand{\coeff}[2]{[#1^{#2}]}
\newcommand{\One}{\mathcal E_M}
\newcommand{\Kcal}{\mathcal K}
\DeclareMathOperator{\Var}{Var}
\hypersetup{pdftitle={Finite-Core Universality and Sharp Large Order for Countable Exponential-Feedback Transseries},pdfauthor={Research draft prepared with ChatGPT},pdfsubject={Sharp cutoff, multiplicative coefficient asymptotics, Gaussian fluctuations, and formal inverse asymptotics}}
\begin{document}
\begin{center}
{\LARGE\bfseries\color{ink} Finite-Core Universality and Sharp Large Order\par for Countable Exponential-Feedback Transseries\par}
\vspace{0.6em}
{\large A sharp cutoff, a moving Gaussian saddle,\par and superquadratic inverse asymptotics\par}
\vspace{0.85em}
Research article prepared for Vladimir Reshetnikov\\
29 September 2026
\end{center}

\begin{abstract}
For the formal equation
\[
 \Uhat(q)=\sum_{j\ge1}q^j\exp\!\bigl(a j^p\Uhat(q)\bigr),
 \qquad a>0,\quad p>1,
\]
we pass from logarithmic large-order estimates to a full multiplicative
coefficient equivalent. The main step is a sharp finite-core reduction.
Keep actions $1,\ldots,M$ in an analytic implicit core, and retain exactly
one further action through its linear response. The coefficients of this
one-tail series are asymptotic to those of $\Uhat$ if and only if
$M(p-1)>1$. At equality the ratio tends to zero, not one. A coupled
small-radius/large-action saddle then gives the full exponential factor
and determinant prefactor using only the finite core. It also yields a
conditional local Gaussian law and an unconditional central limit theorem
for the largest action, centered at an explicitly defined moving saddle,
with standard deviation asymptotic to
$\sqrt{pn}/((p-1)\log n)$.

The result allows nonnegative changes to finitely many initial slopes and
weights, with the first weight kept equal to one. For $p>2$ the finite core disappears from the leading
multiplicative equivalent. For $p=2$ a factor
$\exp((1+1/a)r_n^2)$ survives, where
$r_n(1+r_n)e^{2r_n}=an$; discarding it produces an unbounded relative error.
For $p>2$ we additionally prove a sharp formal-reversion law,
$v_n=-u_n+(n+1)u_2u_{n-1}+o(nu_{n-1})$.
The proofs are self-contained apart from elementary analytic and formal
power-series tools. Exact coefficient checks, numerical saddle diagnostics,
a proof-dependency audit, and a further research program accompany the text.
\end{abstract}

\noindent\textbf{Status.}
This is an AI-assisted research draft with conventional mathematical proofs,
not a Lean-verified development. The model-specific theorems below are the
proposed contributions. Their global originality and publication priority
have not been established. No named longstanding conjecture, general
summability theorem, or full inverse theorem in the quadratic/subquadratic
regime is claimed solved.

\clearpage
\setcounter{tocdepth}{1}\hypersetup{bookmarksdepth=2}
\tableofcontents
\clearpage

\section{The problem left by the repository}
\label{sec:context}

\subsection{A precise extension, rather than another coefficient list}

The ProveIt transseries project develops asymptotic scales, power--logarithmic
and exponential expansions, formal composition and reversion, analytic
realization, and quantitative inversion. Its current documentation carefully
distinguishes the scope of the Lean modules from the human research
manuscripts \cite{repo}. In particular, a formal expansion and a theorem
about the growth or realization of its coefficients are different objects.

The recently filed article \emph{A sharp regularity classification for
countable exponential-feedback transseries} \cite{feedback} studies
\[
 \Uhat(q)=\sum_{j\ge1}q^j e^{\lambda_j\Uhat(q)}.
\]
It obtains convergence and Gevrey classifications, logarithmic coefficient
asymptotics, macroscopic concentration, and sectorial realizations. After its
concentration proof it explicitly leaves a full multiplicative equivalent,
a local limit theorem, and the fluctuation scale around a sufficiently
accurate finite-$n$ center unresolved. Its reversion result transfers Gevrey
class membership, not sharp coefficient asymptotics.

Those are the gaps addressed here, for the substantial but specific class
$\lambda_j=a j^p$, $p>1$, and finite nonnegative perturbations of that class.
The distinction in precision is important. Knowing
$\log u_n=X_n+o(n)$ allows an undetermined factor
$\exp(o(n))$. Such a factor can be larger than every power of $n$.
We exhibit exactly this phenomenon at $p=2$ and give a finite procedure
that incorporates all factors needed for a relative error tending to zero.

The source inspected for \cite{feedback} is pinned at repository commit
\path{5804c7aff9955e6cbc09be2a71a295a07b238bc9}; its TeX blob is
\path{c6a83d71973a12fe7ed970d2abc56e2b49e6808f}.
The relevant current directory is
\begin{quote}\small
\path{Analysis/Transseries/docs/series-and-transseries/Exponential_Feedback_Regularity_Classification/}.
\end{quote}
The repository was changing during inspection. We do not claim an exhaustive
audit of every incoming archive or every branch. Our main proofs do not
assume an unreviewed theorem from an earlier research draft: the needed
coefficient identity and macroscopic bounds are established again below.

\subsection{Relation to established mathematics}

Lagrange inversion is classical; Gessel's survey \cite{gessel} provides several
formal and combinatorial forms. Infinite-colour inversion is also an
established subject, with considerably broader settings than our countable
scalar kernel \cite{jkt}. Neither Lagrange inversion nor the passage to
infinitely many formal colours is claimed as a new method here.

Condensation, finite clouds of small components, and threshold phenomena
already occur in superexponentially weighted random trees. Janson, Jonsson,
and Stef\'ansson \cite{jjs} obtain detailed results for product degree weights,
and Janson \cite{janson} surveys the relation to random allocations. Our
coefficient weights instead contain
$(\lambda_{j_1}+\cdots+\lambda_{j_m})^{m-1}$, not a fixed product of
single-part weights. The growing action has scale $n/\log n$, and the
small-action threshold carries logarithmic factors. Thus the familiar
condensation picture is useful context, but is not itself our novelty claim
or a substitute for the proofs below.

Rapid-coefficient expansions under formal operations have a classical
antecedent in Bender's work \cite{bender}. We give the specific convolution
and reversion argument needed here, including its range $p>2$, rather than
asserting that formal reversion automatically preserves a leading constant.
For the formal transseries viewpoint, see also Edgar \cite{edgar}.

\subsection{What the results do and do not provide}

The main output is a multiplicative asymptotic theorem for a divergent
one-generator transseries. Substituting $q=e^{-x}$ gives
$\sum u_ne^{-nx}$, whose support is entirely elementary; the difficult
information lies in the coefficient growth and the implicit feedback.
No enlarged monomial support is required for the phenomenon.

The finite-core formula uses analytic functions only near the origin, where
they are unambiguous. It never evaluates the original infinite kernel at a
positive argument. This point is essential: for $p>1$, that infinite sum
diverges for every fixed positive pair $(q,u)$. Similarly, no Borel or
Laplace integral is justified merely by obtaining precise large-order
coefficients. Those analytic questions remain separate.

\section{Model, finite core, and principal theorems}
\label{sec:main}

\subsection{The admissible family}

Fix $a>0$, $p>1$, and an integer $M_0\ge1$. Assume
\begin{equation}
 w_1=1,\qquad w_j,\lambda_j\ge0,\qquad
 w_j=1,\quad\lambda_j=a j^p\quad(j>M_0).
 \label{eq:assumptions}
\end{equation}
The finitely many exceptional values are finite real numbers. Write
\begin{equation}
 \Uhat(q)=\sum_{j\ge1}w_jq^j e^{\lambda_j\Uhat(q)}
        =\sum_{n\ge1}u_nq^n.
 \label{eq:model}
\end{equation}
All infinite identities involving $\Uhat$ are formal. Coefficients of any
fixed degree involve only finitely many summands. We prove below that the
solution exists uniquely and $u_1=1$.
The \emph{pure model} has $w_j=1$, $\lambda_j=a j^p$ for all $j\ge1$.

Fix an integer $M\ge M_0$. Define the analytic core and response by
\begin{align}
 \Phi_M(t,u)&=\sum_{j=1}^M w_jt^je^{\lambda_ju},
 &G_M(t)&=\Phi_M(t,G_M(t)),\\
 R_M(t)&=\bigl(1-\partial_u\Phi_M(t,G_M(t))\bigr)^{-1}.
 \label{eq:core}
\end{align}
The branches have $G_M(0)=0$, $G_M'(0)=1$, and $R_M(0)=1$.
They are analytic on a fixed complex neighborhood of zero. Define
\begin{equation}
 W_M(q)=R_M(q)\sum_{j>M}q^j e^{a j^pG_M(q)},
 \qquad w_{n,M}=[q^n]W_M(q).
 \label{eq:one-tail}
\end{equation}
The notation $w_{n,M}$ is distinct from the primitive weights $w_j$.
Although $G_M$ and $R_M$ are analytic, $W_M$ is a \emph{formal} series.
Its coefficient $w_{n,M}$ counts exactly one primitive action above $M$;
it is not the coefficient of the truncation $G_M$ itself.

\begin{theorem}[Sharp finite-core cutoff]
\label{thm:cutoff}
Under \eqref{eq:assumptions},
\begin{equation}
 \boxed{\quad
 \lim_{n\to\infty}\frac{w_{n,M}}{u_n}=
 \begin{cases}
 1,&M(p-1)>1,\\[2pt]
 0,&M(p-1)\le1.
 \end{cases}\quad}
 \label{eq:cutoff}
\end{equation}
For the pure model the smallest sufficient integer is
\begin{equation}
 M_*(p)=\left\lfloor\frac1{p-1}\right\rfloor+1.
 \label{eq:mincore}
\end{equation}
All parameters, including $M$, are fixed as $n\to\infty$.
\end{theorem}

% ed. (2026-09-29): editorial cross-reference to the later Poisson-layer package (batch 48), which sharpens the boundary case
\begin{ednote}
The boundary case $M(p-1)=1$ is sharpened, consistently with
\eqref{eq:cutoff}, by the later package \cite{ed:plb}. At $p=1+1/M$ its
Theorem~\texttt{thm:boundary} gives
$u_n\sim w_{n,M}\exp(r_n^{M+1}/a^M)$, an explicit missing factor that
makes the ratio tend to zero as stated here; for $M=1$, $p=2$ that factor
is $\exp(r_n^2/a)$, a part of the factor in \eqref{eq:quadratic}. A finite
limit $\ee^{-\lambda(\theta)}$ of the ratio, with a Poisson law for the
first omitted actions, occurs there only in a window where the exponent
$p=p_n$ approaches $1+1/M$ with $n$ (Theorem~\texttt{thm:window}), outside
the fixed-parameter setting of this theorem.
\end{ednote}

\subsection{A computable multiplicative equivalent}

In the next theorem let $G=G_M$, $R=R_M$ to simplify notation.
The saddle is a real pair $(z_n,\tau_n)$, with $z_n$ not required to be an
integer. Set $k_n=n-z_n$ and $\Lambda_n=a z_n^p$.

\begin{theorem}[Finite-core saddle formula]
\label{thm:saddle}
For every fixed $M\ge M_0$, there is, for all sufficiently large $n$, a
unique solution in the asymptotic region specified below to
\begin{equation}
 \boxed{\qquad
 n=z_n+a z_n^p\tau_nG'(\tau_n),\qquad
 -\log\tau_n=apz_n^{p-1}G(\tau_n).
 \qquad}
 \label{eq:saddle}
\end{equation}
It satisfies
\begin{align}
 z_n&\sim\frac{p}{p-1}\frac{n}{\log n},
 &\tau_n&\sim \frac1a\left(\frac{p-1}{p}\right)^p
                   \frac{(\log n)^p}{n^{p-1}}.
 \label{eq:saddlescales}
\end{align}
Define, with all core functions evaluated at $\tau_n$,
\begin{align}
 B_n&=\Lambda_n(\tau_nG'+\tau_n^2G''),
 &C_n&=1+\frac{pk_n}{z_n},\\
 H_n&=ap(p-1)z_n^{p-2}G,
 &\Delta_n&=C_n^2-H_nB_n.
 \label{eq:determinant}
\end{align}
Then $\Delta_n>0$ eventually and
\begin{equation}
 \boxed{\quad
 w_{n,M}\sim\Acal_{n,M}:=
 \frac{R_M(\tau_n)}{\sqrt{\Delta_n}}
 \exp\!\left(a z_n^pG_M(\tau_n)-(n-z_n)\log\tau_n\right).
 \quad}
 \label{eq:full-equivalent}
\end{equation}
Consequently $u_n\sim\Acal_{n,M}$ when $M(p-1)>1$, whereas
$\Acal_{n,M}/u_n\to0$ when $M(p-1)\le1$.
\end{theorem}

The uniqueness assertion is local to, for example,
$z\log n/n\in[c_-,c_+]$, where
$0<c_-<p/(p-1)<c_+$ are fixed and sufficiently close, and to the
small branch of the core. It is not a claim about every possible positive
solution of the displayed equations away from that region.

\subsection{Fluctuations and inverse asymptotics}

The positive Lagrange formula in Section~\ref{sec:formal} defines a
probability distribution on ordered action configurations of total action
$n$. Let $J_n$ denote their largest action and let $\One$ be the event
that exactly one action exceeds $M$.

\begin{theorem}[Moving Gaussian law]
\label{thm:clt}
For every fixed $M\ge M_0$, put
\begin{equation}
 \sigma_n^2=\frac{B_n}{\Delta_n}
       \sim\frac{z_n^2}{pn}
       \sim\frac{p}{(p-1)^2}\frac{n}{(\log n)^2}.
 \label{eq:variance}
\end{equation}
For each fixed $A<\infty$, uniformly over integers
$|j-z_n|\le A\sigma_n$,
\begin{equation}
 \sigma_n\Prob(J_n=j\mid\One)
 =\frac1{\sqrt{2\pi}}
   \exp\!\left(-\frac{(j-z_n)^2}{2\sigma_n^2}\right)+o(1).
 \label{eq:local-clt}
\end{equation}
The conditional standardized variables converge in distribution to
$\Normal(0,1)$. If $M(p-1)>1$, then also
\begin{equation}
 \frac{J_n-z_n}{\sigma_n}\ \Longrightarrow\ \Normal(0,1)
 \quad\text{without conditioning}.
 \label{eq:clt}
\end{equation}
\end{theorem}

We deliberately distinguish the conditional \emph{local} theorem from the
unconditional \emph{weak} theorem. A conditioning event of probability
$1-o(1)$ transfers weak convergence, but alone does not transfer errors on
the smaller $1/\sigma_n$ scale of individual point probabilities.

For the pure model, let $r_n>0$ be the unique solution of
\begin{equation}
 r_n(1+r_n)^{p-1}e^{pr_n}=a n^{p-1},
 \qquad j_n^0=\frac n{1+r_n},\quad k_n^0=n-j_n^0,
 \label{eq:bare-r}
\end{equation}
and define
\begin{equation}
 S_n=\frac{\exp\!\bigl(k_n^0(1+pr_n)\bigr)}
              {\sqrt{1+2pr_n+pr_n^2}}.
 \label{eq:bare-S}
\end{equation}

\begin{theorem}[Explicit large order and superquadratic reversion]
\label{thm:explicit}
For the pure model, if $p>3/2$ then
\begin{equation}
 u_n\sim S_n
  \exp\!\left(\frac{(a+1)(k_n^0)^2}{a(j_n^0)^p}\right).
 \label{eq:first-correction}
\end{equation}
In particular, $u_n\sim S_n$ for $p>2$, whereas for $p=2$,
\begin{equation}
 \boxed{\qquad u_n\sim S_n\exp\!\bigl((1+1/a)r_n^2\bigr).\qquad}
 \label{eq:quadratic}
\end{equation}
If $p>2$ and $\Vhat=\Uhat^{\langle-1\rangle}=\sum v_nq^n$, then
\begin{equation}
 v_n=-u_n+(n+1)u_2u_{n-1}+o(nu_{n-1}),
 \qquad v_n\sim-u_n\sim-S_n.
 \label{eq:inverse}
\end{equation}
Moreover,
\begin{equation}
 \frac{u_n}{u_{n-1}}\sim e^{pr_n}
 =\frac{a n^{p-1}}{r_n(1+r_n)^{p-1}}.
 \label{eq:ratio}
\end{equation}
\end{theorem}

% ed. (2026-09-29): editorial cross-reference to the microscopic-condensation package (batch 48), an independent derivation of eq:quadratic at a=1
\begin{ednote}
For $a=1$ and slopes eventually equal to $j^2$, \eqref{eq:quadratic}
(with Proposition~\ref{prop:first-correction} for finitely many changed
slopes) is also the main coefficient theorem, Theorem~\texttt{thm:coefficient},
of the package \cite{ed:mic}, derived independently and without citing this
article: the same root $r_n(1+r_n)\ee^{2r_n}=n$, the same exponent and the
same prefactor. That article adds slope tails $\alpha j+\eta j/\log j$,
the rate $O((\log n)^2/\sqrt n)$, a total-variation Poisson law for the
number of actions equal to two jointly with the largest action, and a
positive-ray Borel equivalent.
\end{ednote}

\section{Exact formal identities and their probability model}
\label{sec:formal}

\begin{lemma}[Formal existence and positive coefficient formula]
\label{lem:lagrange}
Equation~\eqref{eq:model} has a unique solution in $q\R[[q]]$, with
$u_1=1$ and nonnegative coefficients. Its coefficients are
\begin{equation}
 u_n=\sum_{m=1}^n\frac1{m!}
 \sum_{\substack{j_1+\cdots+j_m=n\\ j_i\ge1}}
 \left(\prod_{i=1}^m w_{j_i}\right)
 \left(\sum_{i=1}^m\lambda_{j_i}\right)^{m-1}.
 \label{eq:positive}
\end{equation}
Equivalently, for finitely supported multiplicity vectors
$\mathbf m=(m_1,m_2,\ldots)$,
\begin{equation}
 u_n=\sum_{\sum j m_j=n}
 \frac{\prod_j w_j^{m_j}}{\prod_jm_j!}
 \left(\sum_j\lambda_jm_j\right)^{\sum_jm_j-1}.
 \label{eq:partition}
\end{equation}
For the one-part term the zeroth power is interpreted as $1$, including
when its base vanishes.
\end{lemma}

\begin{proof}
The coefficient of $q^n$ on the right side of \eqref{eq:model} uses only
coefficients of $\Uhat$ of degrees at most $n-1$. This recursively defines
a unique solution and proves nonnegativity and $u_1=w_1=1$.

For the explicit expression, introduce a separate marker $s$ and solve
$U_s=s\Phi(q,U_s)$, where
$\Phi(q,u)=\sum_{j\ge1}w_jq^je^{\lambda_ju}$.
Classical formal Lagrange inversion gives
\[
 [s^m]U_s=\frac1m[u^{m-1}]\Phi(q,u)^m.
\]
Expanding the product and the exponential gives \eqref{eq:positive} after
setting $s=1$. That substitution is valid coefficientwise: every factor
of $\Phi$ has positive $q$-order, so only $m\le n$ can contribute to
$[q^n]$. Grouping ordered lists by their multiplicities gives
\eqref{eq:partition}.
\end{proof}

For all sufficiently large $n$, the one-part configuration of size $n$
has weight $1$, so $u_n>0$. Normalize every ordered summand in
\eqref{eq:positive} by $u_n$. This defines a probability measure on the
finite disjoint union of all compositions of $n$. Write $K_n=m$ for its
number of parts, $J_n=\max_i j_i$, and
\begin{equation}
 E_n=n-K_n+1=1+\sum_i(j_i-1)
 \label{eq:excess}
\end{equation}
for one plus the total excess above unit actions. The integer $E_n$ is
not the largest part, although the two become asymptotically close.

\begin{lemma}[The exact one-tail identity]
\label{lem:marking}
The core functions in \eqref{eq:core} are analytic near zero with
nonnegative Taylor coefficients. Furthermore,
\begin{equation}
 w_{n,M}=u_n\Prob(\One).
 \label{eq:prob-one}
\end{equation}
\end{lemma}

\begin{proof}
The analytic implicit function theorem applies to
$u-\Phi_M(t,u)$ at $(0,0)$, since its derivative in $u$ is $1$ there.
Formal iteration proves nonnegative coefficients, and the analytic
solution has the same Taylor series. The series
$P(t)=\partial_u\Phi_M(t,G_M(t))$ has nonnegative coefficients and zero
constant term; hence $R_M=(1-P)^{-1}$ does too.

Mark every primitive action $j>M$ by $s$:
\[
 U_s=\Phi_M(q,U_s)+s\sum_{j>M}q^je^{a j^pU_s}.
\]
At $s=0$ the solution is $G_M$. Differentiating formally at $s=0$ yields
\[
 (1-\partial_u\Phi_M(q,G_M))\left.\partial_sU_s\right|_{s=0}
 =\sum_{j>M}q^je^{a j^pG_M}.
\]
Thus $[s]U_s=W_M$. In \eqref{eq:positive}, the power of $s$ is exactly
the number of parts above $M$, proving \eqref{eq:prob-one}.
\end{proof}

\begin{example}[Small coefficients in the pure model]
Directly from \eqref{eq:partition},
\begin{align*}
 u_1&=1,\qquad u_2=a+1,\\
 u_3&=1+a(1+2^p)+\tfrac32a^2,\\
 u_4&=1+a(1+2^p+3^p)
       +\tfrac12a^2(2+2^p)^2+\tfrac83a^3.
\end{align*}
For $a=1,p=2$, $u_3=15/2$ and the inverse coefficient is
$v_3=2u_2^2-u_3=1/2>0$. Thus an eventual-sign theorem for the inverse
must not be misread as a sign assertion at every small degree.
\end{example}

\section{Macroscopic localization, with full estimates}
\label{sec:macro}

The following lemma establishes the coarse bounds needed later. The
logarithmic conclusion is included to make the reduction self-contained;
it is not presented as new relative to \cite{feedback}.
Put $s=p-1$ and
\begin{equation}
 X_n=s n\log n-p n\log\log n
 +\left(\log a+p\log\frac p{s}-s\right)n.
 \label{eq:Xn}
\end{equation}

\begin{lemma}[Exponential-scale localization]
\label{lem:macro}
Under \eqref{eq:assumptions}, $\log u_n=X_n+o(n)$. For every
$\delta>0$ there is $c_\delta>0$ such that
\begin{equation}
 \Prob\!\left(\left|\frac{E_n\log n}{n}-\frac p{p-1}\right|>\delta\right)
 \le\exp(-c_\delta n+o(n)).
 \label{eq:macro-excess}
\end{equation}
For every $\varepsilon>0$ there is $d_\varepsilon>0$ such that
\begin{equation}
 \Prob\bigl(J_n<(1-\varepsilon)E_n\bigr)
 \le\exp(-d_\varepsilon n+o(n)).
 \label{eq:macro-giant}
\end{equation}
Also, for every fixed $M\ge M_0$,
$\log w_{n,M}=X_n+o(n)$, even if $M$ is below the sharp cutoff.
\end{lemma}

\begin{proof}
Choose constants $D\ge0$ and $W\ge1$ so that
$\lambda_j\le a j^p+D$ and $w_j\le W$ for all $j$. In a composition
with $m$ parts put $r=n-m+1$. Convexity of $x^p$ and repeated transfer
of excess to a largest part give
\[
 \sum_i j_i^p\le r^p+m-1.
\]
There are at most $r-1$ nonunit parts, so their primitive weight product
is at most $W^{r-1}$. Increasing $D$ absorbs the extra $a(m-1)$ and gives
the row bound
\begin{equation}
 [\text{$m$-part contribution}]
 \le\binom{n-1}{r-1}W^{r-1}
       \frac{(a r^p+Dm)^{m-1}}{m!}.
 \label{eq:row-upper}
\end{equation}
Conversely, the configuration with one action $r$ and $n-r$ unit actions
contributes at least
\begin{equation}
 T_{n,r}:=\frac{(a r^p)^{n-r}}{(n-r)!},
 \qquad r>M_0.
 \label{eq:lower-spike}
\end{equation}
The factor counting its possible positions cancels the difference between
$m!$ and $(m-1)!$.

Let $L=\log n$. Taking $r=\lfloor(p/s)n/L\rfloor$ in
\eqref{eq:lower-spike} and using Stirling's formula gives
$\log T_{n,r}=X_n+o(n)$.
We now show why no row improves this logarithm by order $n$.
If $r\ge\eta n$ for fixed $\eta>0$, the logarithm of the right side of
\eqref{eq:row-upper} is at most
$(1-\eta)s nL+O(n\log L)$, using the elementary entropy bound
$\binom{n-1}{r-1}\le2^n$. Such rows are excluded from any maximizing
sequence, so necessarily $r/n\to0$.

In that range the binomial entropy and $(r-1)\log W$ are $o(n)$.
Comparison with the lower bound forces
$\log(a r^p+Dm)\ge pL-O(\log L)$; hence
$\log r=L-O(\log L)$ and $Dm=o(r^p)$. Write $\alpha=r/n$.
Uniform Stirling estimates now give, for any row within $O(n)$ of the
maximum,
\[
 \frac1n\log[\text{row bound}]
 =(1-\alpha)\{sL+p\log\alpha+\log a
                -\log(1-\alpha)+1\}+o(1).
\]
Comparison once more with $X_n/n$ gives $\alpha L=O(\log L)$.
Consequently $\alpha\log\alpha=o(1)$ and, putting $y=rL/n$,
\begin{equation}
 \frac1n\log[\text{row bound}]
 =sL-p\log L+\log a+1+p\log y-sy+o(1).
 \label{eq:variational}
\end{equation}
The function $g(y)=p\log y-sy$ tends to $-\infty$ at both ends of
$(0,\infty)$ and is strictly concave, with its unique maximum at $p/s$.
This proves the upper logarithmic estimate. It also proves uniform
exponential loss away from any fixed neighborhood of $p/s$: otherwise a
sequence of offending rows would, by the preceding localization, have a
limit point contradicting the strict maximum of $g$. Summing at most $n$
rows changes a logarithm by only $O(\log n)$. This proves
\eqref{eq:macro-excess}.

For \eqref{eq:macro-giant}, it is enough to take $0<\varepsilon<1/2$.
Restrict $r$ to a fixed compact multiple of $n/L$, since its complement
already has exponentially small weight. Under the constraint
$j_i\le b=\lfloor(1-\varepsilon)r\rfloor$, convexity maximizes the slope
sum by putting the excess into at most two nonunit parts, of sizes $b$
and $r-b+1$. Indeed, transferring mass between two interior nonunit parts
increases the convex sum until one reaches $1$ or the upper cap.
Thus the slope sum is at most
\[
 a\{b^p+(r-b+1)^p\}+D'm
 =a r^p\{(1-\varepsilon)^p+\varepsilon^p+o(1)\}.
\]
The expression in braces tends to a constant strictly below $1$. Raising
to the power $m-1\sim n$ loses a fixed exponential factor. The
composition entropy and primitive weights are only $\exp(o(n))$ in this
range and cannot compensate. The same maximum calculation proves
\eqref{eq:macro-giant}.

Finally, $w_{n,M}\le u_n$, while the lower configuration in
\eqref{eq:lower-spike} has exactly one part above any fixed $M$ for large
$n$. Therefore it is also a lower bound for $w_{n,M}$, proving its
logarithmic asymptotic.
\end{proof}

A useful consequence is that, with exponentially high probability,
$K_n\sim n$, $E_n\asymp n/\log n$, and the largest action is unique.
This still does \emph{not} imply that all remaining actions are bounded
by a particular $M$. That stronger assertion is the next step.

\section{Why a sufficient finite core captures almost all weight}
\label{sec:cutoff-proof}

\begin{proposition}[Sufficiency of the strict cutoff]
\label{prop:sufficiency}
If $M(p-1)>1$, then $\Prob(\One)\to1$.
\end{proposition}

\begin{proof}
Put $c_*=p/(p-1)$. Fix a small width $\delta>0$ and a small
$\varepsilon\in(0,1/3)$, to be chosen below. By
Lemma~\ref{lem:macro}, except on a set of probability
$\exp(-cn+o(n))$, a configuration has
\begin{equation}
 (c_*-\delta)\frac n{\log n}\le r=E_n
 \le(c_*+\delta)\frac n{\log n},
 \qquad L_0=J_n\ge(1-\varepsilon)r.
 \label{eq:good-region}
\end{equation}
The largest part $L_0$ is unique for large $n$, and is above $M$.
Every other part $i>1$ has $i-1\le\varepsilon r$.

Consider a good configuration with an additional part $i>M$. Choose one
such part by a fixed deterministic rule, for example its earliest position.
Replace the pair $(L_0,i)$ by $(L_0+i-1,1)$, keeping all positions and the
number $m$ of parts unchanged. Both nonunit parts in this operation are
in the pure tail, while the new unit part has weight $w_1=1$. Thus the
primitive weight product and the denominator $m!$ in
\eqref{eq:positive} are unchanged.

Let $S$ and $S'$ be the old and new slope sums. For every $\eta>0$,
choosing $\varepsilon$ small enough gives, uniformly in
\eqref{eq:good-region} and in $i>M$,
\begin{equation}
 S'-S\ge a(p-\eta)r^{p-1}(i-1).
 \label{eq:gain}
\end{equation}
Indeed, the mean-value theorem bounds the gain at the large part below by
$ap(1-\varepsilon)^{p-1}r^{p-1}(i-1)$. The loss is
$a i^p-\lambda_1$. Since $\lambda_1$ is fixed and nonnegative, its possible
contribution can be discarded, and
$i^p/(i-1)\le C_p(\varepsilon r+1)^{p-1}$ for $i\ge2$.
Making $\varepsilon$ small absorbs this term and proves
\eqref{eq:gain} for large $n$.

By the convex bound used earlier,
$S'\le a r^p+Dm=a r^p(1+o(1))$ uniformly here. The old coefficient
weight divided by the new one is therefore at most
\begin{equation}
 \left(\frac S{S'}\right)^{m-1}
 \le\exp\!\left(-\frac{(p-2\eta)(m-1)(i-1)}r\right)
 \label{eq:weight-ratio}
\end{equation}
for sufficiently large $n$.

We record the multiplicity of this map, rather than treating it as an
unjustified injection. In the image, the unique giant is still recognizable.
Given its image configuration, the chosen slot, and $i$, the original
configuration is recovered by subtracting $i-1$ from the giant and replacing
the unit in that slot by $i$. Thus there are at most $m$ preimages for each
integer $i$. Let $\mathcal B_n$ denote the good configurations with an
additional part above $M$. Summing \eqref{eq:weight-ratio} gives
\begin{align}
 \Prob(\mathcal B_n)
 &\le \sup_{r\text{ in }\eqref{eq:good-region}}
 n\sum_{i=M+1}^{\infty}
  \exp\!\left(-\frac{(p-2\eta)(n-r)(i-1)}r\right)\notag\\
 &\le n^{\,1-\frac{(p-2\eta)M}{c_*+\delta}+o(1)}.
 \label{eq:small-tail-bound}
\end{align}
The denominator of the geometric series tends to $1$ and does not change
the power of $n$. Because $M(p-1)>1$, we may choose $\eta$ and $\delta$
small enough that the exponent in the last line is negative.
The exceptional configurations have exponentially small probability.
On the good event at least one part exceeds $M$, so the absence of a second
one is precisely $\One$. This proves the proposition.
\end{proof}

\begin{remark}[What makes this stronger than macroscopic condensation]
The large-action result only says that the remaining excess is $o(n/\log n)$.
The mass-transfer estimate controls every extra action above a \emph{fixed}
cutoff, and pays explicitly for its possible locations. The competition
between a factor of order $n$ and a suppression of order
$n^{-M(p-1)}$ is the origin of the cutoff. It also explains why an
unqualified statement that ``one large action dominates'' does not establish
a multiplicative coefficient formula.
\end{remark}

\section{A uniform small-radius coefficient saddle}
\label{sec:inner-saddle}

We next prove the analytic estimate used to compute $w_{n,M}$.
Only a fixed analytic core is involved. In particular, no exchange of an
infinite divergent tail with a contour integral will be made.

\begin{lemma}[Coefficient estimate with a shrinking saddle]
\label{lem:inner}
Let $G(t)=t+O(t^2)$ and $R(t)=1+O(t)$ be analytic at zero with
nonnegative Taylor coefficients. Fix $a>0,p>1$, and constants
$0<c_-<c_+<\infty$. For integers
\[
 c_-\frac n{\log n}\le j\le c_+\frac n{\log n},\qquad
 k=n-j,\quad\Lambda=a j^p,
\]
let $t=t_j>0$ solve
\begin{equation}
 \Lambda tG'(t)=k.
 \label{eq:inner-saddle}
\end{equation}
Then this solution exists uniquely on the small positive branch for large
$n$, $t\sim k/\Lambda\to0$ uniformly, and
\begin{equation}
 [q^k]R(q)e^{\Lambda G(q)}
 =\frac{R(t)e^{\Lambda G(t)}t^{-k}}
        {\sqrt{2\pi B_j}}(1+o(1)),
 \qquad B_j=\Lambda(tG'(t)+t^2G''(t)).
 \label{eq:inner-estimate}
\end{equation}
The $o(1)$ is uniform over the indicated integers $j$.
\end{lemma}

\begin{proof}
Since $tG'(t)=t+O(t^2)$ has positive derivative near zero, the small solution
exists uniquely whenever $k/\Lambda$ is sufficiently small. In the stated
range,
$k/\Lambda\asymp n^{1-p}(\log n)^p\to0$ uniformly. Also
$B_j\sim k\sim n$.

Cauchy's formula on $q=te^{i\theta}$ gives
\begin{align*}
 [q^k]R(q)e^{\Lambda G(q)}
 =\frac{e^{\Lambda G(t)}t^{-k}}{2\pi}
 \int_{-\pi}^{\pi}R(te^{i\theta})
 \exp\!\bigl(\Lambda[G(te^{i\theta})-G(t)]-ik\theta\bigr)\,\dd\theta.
\end{align*}
The linear term vanishes by \eqref{eq:inner-saddle}. With
$D=t\,\dd/\dd t$, analyticity and $G(t)=t+O(t^2)$ give
$\Lambda D^hG(t)=O(k)$ for every fixed positive integer $h$, uniformly
as $t\to0$. Hence, for $|\theta|\le k^{-2/5}$,
\[
 \Lambda[G(te^{i\theta})-G(t)]-ik\theta
 =-\tfrac12B_j\theta^2+O(k|\theta|^3),
\]
where the error is $o(1)$ throughout this arc. The response satisfies
$R(te^{i\theta})/R(t)=1+o(1)$ there.

For the complement, write $G(q)=\sum_{h\ge1}g_hq^h$, with $g_1=1$
and $g_h\ge0$. Then
\begin{equation}
 G(t)-\Re G(te^{i\theta})
 =\sum_{h\ge1}g_ht^h(1-\cos h\theta)
 \ge t(1-\cos\theta).
 \label{eq:circle-decay}
\end{equation}
Also $|R(te^{i\theta})|\le R(t)$. Since $\Lambda t\asymp k$, the
complementary integral is bounded by the main positive scale times
$O(e^{-c k^{1/5}})$, using $1-\cos\theta\ge c\theta^2$ on
$[-\pi,\pi]$. The central Gaussian integral is
$\sqrt{2\pi/B_j}(1+o(1))$. This proves
\eqref{eq:inner-estimate} and its uniformity.
\end{proof}

The same proof allows a fixed analytic multiplier $t^d h(t)$ with
$h(0)>0$: its coefficient contribution is asymptotic to its value at the
saddle times the right side of \eqref{eq:inner-estimate}. The factor
$e^{id\theta}$ is $1+o(1)$ on the central arc. We will use this observation
with fixed $d=M+1$ in the sharpness proof.

\section{The moving action saddle and its determinant}
\label{sec:outer-saddle}

\subsection{The exponent after coefficient extraction}

For real $j$ in a fixed compact multiple of $n/\log n$, define $t_j$ by
\eqref{eq:inner-saddle} and set
\begin{equation}
 f_n(j)=a j^pG(t_j)-(n-j)\log t_j.
 \label{eq:outer-exponent}
\end{equation}
Here and throughout this section $G=G_M$, $R=R_M$.
Differentiating the defining equation for $t_j$ gives
\begin{equation}
 \frac{t_j'}{t_j}=-\frac{C_j}{B_j},
 \qquad C_j=1+\frac{p(n-j)}j.
 \label{eq:t-derivative}
\end{equation}
The saddle condition eliminates the terms containing $t_j'$ when
$f_n$ is differentiated once, so
\begin{align}
 f_n'(j)&=apj^{p-1}G(t_j)+\log t_j,\\
 f_n''(j)&=ap(p-1)j^{p-2}G(t_j)-\frac{C_j^2}{B_j}.
 \label{eq:f-derivatives}
\end{align}
These are exact identities, not leading-order approximations.

\begin{lemma}[Location, curvature, and uniqueness]
\label{lem:outer}
Choose fixed $c_-<p/(p-1)<c_+$. Uniformly on the interval
$I_n=[c_-n/\log n,c_+n/\log n]$,
\begin{align}
 \frac{f_n'(c n/\log n)}{\log n}
 &=\frac pc-(p-1)+o(1),\\
 f_n''(j)&=-\frac{pn}{j^2}(1+o(1)),
 & f_n'''(j)&=O\!\left(\frac{(\log n)^3}{n^2}\right).
 \label{eq:curvature}
\end{align}
For large $n$ there is therefore exactly one zero $z_n$ of $f_n'$ in
$I_n$, and it is the strict maximum of $f_n$ there.
The pair $(z_n,t_{z_n})$ is the solution asserted in
Theorem~\ref{thm:saddle}.
\end{lemma}

\begin{proof}
In this interval $k=n-j\sim n$, $G(t)=t(1+O(t))$,
$G'(t)=1+O(t)$, and $t\sim k/(a j^p)$. Thus
\[
 apj^{p-1}G(t_j)=\frac{pk}{j}(1+O(t_j)),
\]
while
\[
 \log t_j=(1-p)\log n+p\log\log n-\log a-p\log c+o(1)
\]
when $j=c n/\log n$. This proves the first estimate and the signs at
the two endpoints. In \eqref{eq:f-derivatives},
\[
 B_j=k(1+O(t_j)),\qquad
 ap(p-1)j^{p-2}G(t_j)=\frac{p(p-1)k}{j^2}(1+O(t_j)).
\]
Since $C_j=1+pk/j$, their difference is
$-pk/j^2-2p/j-1/k+o(n/j^2)$, proving the curvature estimate.
Differentiating these analytic expressions, using
$t_j'/t_j=O(1/j)$ and the corresponding derivatives of
$G(t)=t+O(t^2)$, gives $f_n'''=O(n/j^3)$, the stated bound.
Strict negative curvature and the endpoint signs prove uniqueness.
The leading location and the leading value of $t_{z_n}$ follow immediately,
giving \eqref{eq:saddlescales}.
\end{proof}

\subsection{Completing the two-stage asymptotic}

\begin{proof}[Proof of Theorem~\ref{thm:saddle}]
First localize the sum in \eqref{eq:one-tail} to $j\in I_n$, choosing
$c_-,c_+$ on opposite sides of $p/(p-1)$. Under the one-tail event,
$j$ is the largest part. Lemma~\ref{lem:macro} shows that the contribution
outside such an interval is at most
$u_n\exp(-cn+o(n))$. Both $u_n$ and $w_{n,M}$ have logarithm
$X_n+o(n)$, so this is also exponentially negligible relative to $w_{n,M}$.
This argument does not require the sufficient-cutoff result.

Inside $I_n$, Lemma~\ref{lem:inner} gives
\begin{equation}
 w_{n,M}\sim
 \sum_{j\in I_n\cap\N}
       \frac{R(t_j)}{\sqrt{2\pi B_j}}e^{f_n(j)}.
 \label{eq:outer-sum}
\end{equation}
At $z_n$, the identities \eqref{eq:f-derivatives} give
\[
 -f_n''(z_n)=\frac{\Delta_n}{B_n}>0.
\]
The estimates above yield
\begin{equation}
 B_n\sim n,\qquad
 \Delta_n\sim\frac{pn^2}{z_n^2},\qquad
 \sigma_n=\sqrt{\frac{B_n}{\Delta_n}}\longrightarrow\infty.
 \label{eq:BD-asymptotics}
\end{equation}
For $h=j-z_n$ with $|h|\le\sigma_n\log n$, Taylor's formula gives
\begin{equation}
 f_n(z_n+h)=f_n(z_n)-\frac{h^2}{2\sigma_n^2}+o(1)
 \label{eq:quadratic-window}
\end{equation}
uniformly: the cubic error is
$O((\log n)^3n^{-2}(\sigma_n\log n)^3)
=O((\log n)^3/\sqrt n)=o(1)$.
In the same window $R(t_j)/\sqrt{B_j}$ varies by a relative $o(1)$,
since its logarithmic derivative is $O(1/j)$.

The curvature bound in \eqref{eq:curvature} suppresses the remaining
terms of $I_n$ by $\exp(-c(\log n)^2)$ after the central window is
removed. A polynomial number of terms cannot repair this loss. Finally,
since $\sigma_n\to\infty$, the Gaussian lattice sum is a Riemann sum:
\[
 \sum_{j\in\mathbb Z}
 e^{-(j-z_n)^2/(2\sigma_n^2)}
 \sim\sqrt{2\pi}\,\sigma_n.
\]
Multiplying this with the prefactor in \eqref{eq:outer-sum} gives
\[
 w_{n,M}\sim
 R(\tau_n)e^{f_n(z_n)}\frac{\sigma_n}{\sqrt{B_n}}
 =\frac{R(\tau_n)e^{f_n(z_n)}}{\sqrt{\Delta_n}},
\]
which is \eqref{eq:full-equivalent}. The positive case of
Theorem~\ref{thm:cutoff}, already proved, transfers the formula to $u_n$
when $M(p-1)>1$. The negative case will follow from the next section.
\end{proof}

\begin{remark}[Why there is no missing factor of $2\pi$]
The coefficient saddle supplies $(2\pi B_n)^{-1/2}$. Summing the action
saddle supplies $(2\pi)^{1/2}\sigma_n$. Their product is exactly
$\Delta_n^{-1/2}$. Using only one saddle would miss both this cancellation
and the moving action scale.
\end{remark}

\section{Sharpness: the boundary core is not sufficient}
\label{sec:necessity}

\begin{proposition}[A single omitted small action forces failure]
\label{prop:necessity}
For any fixed $M\ge M_0$, let $d=M+1$. There is a nonnegative subseries
of \eqref{eq:model}, disjoint from the one-tail subseries, whose coefficient
$b_{n,M}$ satisfies
\begin{equation}
 \frac{b_{n,M}}{w_{n,M}}\sim k_n\tau_n^M.
 \label{eq:two-tail-ratio}
\end{equation}
In particular, if $M(p-1)\le1$, then $w_{n,M}/u_n\to0$.
\end{proposition}

\begin{proof}
Temporarily add the action $d$ to the core with a marker $s$:
\begin{equation}
 G_s(q)=\Phi_M(q,G_s(q))+s q^d e^{a d^pG_s(q)}.
 \label{eq:marked-small}
\end{equation}
Set
\[
 R_s(q)=\bigl(1-\partial_u\Phi_M(q,G_s(q))
                 -s a d^p q^d e^{a d^pG_s(q)}\bigr)^{-1}.
\]
At $s=0$, $G_s=G$ and $R_s=R$. Differentiation gives
\begin{align}
 D(q):=\left.\partial_sG_s(q)\right|_{s=0}
 &=R(q)q^d e^{a d^pG(q)}=q^d(1+O(q)),\\
 \dot R(q):=\left.\partial_sR_s(q)\right|_{s=0}
 &=R(q)^2\left(\partial_u^2\Phi_M(q,G(q))D(q)
                       +a d^p q^d e^{a d^pG(q)}\right).
 \label{eq:marked-derivatives}
\end{align}
Both series have nonnegative coefficients, and $\dot R(q)=O(q^d)$.

Now retain exactly one action $j>d$, and differentiate its response
with respect to $s$ at zero. The result is the formal nonnegative series
\begin{equation}
 B_M(q)=\sum_{j>d}q^j e^{a j^pG(q)}
             \left(\dot R(q)+a j^pR(q)D(q)\right).
 \label{eq:two-tail-series}
\end{equation}
By the two independent markers, it counts configurations with exactly one
part $j>d$, exactly one part $d$, and all remaining parts at most $M$.
Consequently it is disjoint from $W_M$, and
$u_n\ge w_{n,M}+b_{n,M}$ for $b_{n,M}=[q^n]B_M$.

For each $j$ in the saddle region, apply Lemma~\ref{lem:inner} and its
analytic-multiplier extension. Relative to the coefficient of
$q^j R(q)e^{a j^pG(q)}$, the second summand in the parentheses in
\eqref{eq:two-tail-series} contributes asymptotically
$a j^p D(t_j)\sim a j^p t_j^d$. The first contributes $O(t_j^d)$ and
is negligible, since $a j^p\to\infty$.
The multiplier $a j^p t_j^d$ varies by a relative $o(1)$ across the
Gaussian action window. Coarse tails are negligible by Lemma~\ref{lem:macro},
since $B_M$ is a subseries of $\Uhat$ and $\log w_{n,M}=X_n+o(n)$.
Inside the coarse interval but outside the Gaussian window, the multiplier
has at most a polynomial effect,
which cannot overcome the saddle decay. The finitely many excluded values
$j\le d$ have merely exponential coefficient growth and are negligible.
Repeating the outer-saddle summation thus gives
\[
 \frac{b_{n,M}}{w_{n,M}}
 \sim a z_n^p\tau_n^d
 =\frac{k_n\tau_n^{d-1}}{G'(\tau_n)}
 \sim k_n\tau_n^M.
\]

From \eqref{eq:saddlescales}, with
$c_q=a^{-1}((p-1)/p)^p>0$,
\begin{equation}
 k_n\tau_n^M
 \sim c_q^M n^{1-M(p-1)}(\log n)^{pM}.
 \label{eq:cutoff-scale}
\end{equation}
This diverges when $M(p-1)<1$, and also at equality, where the positive
power of $\log n$ remains. Since
$u_n/w_{n,M}\ge1+b_{n,M}/w_{n,M}$, the claimed failure follows.
\end{proof}

Proposition~\ref{prop:necessity} completes Theorem~\ref{thm:cutoff} and
the negative-cutoff conclusion of Theorem~\ref{thm:saddle}. It also
identifies exactly what is wrong with replacing the strict inequality in
\eqref{eq:mincore} by a non-strict one. The boundary does not have a
finite nonzero missing-action correction; that correction is already
unbounded.

\section{Gaussian fluctuations around the correct moving center}
\label{sec:fluctuations}

\begin{proof}[Proof of Theorem~\ref{thm:clt}]
On $\One$, the unique action $j>M$ is the largest. Its exact probability
is
\begin{equation}
 \Prob(J_n=j\mid\One)
 =\frac{[q^{n-j}]R_M(q)e^{a j^pG_M(q)}}{w_{n,M}},\qquad j>M.
 \label{eq:conditional-law}
\end{equation}
Use \eqref{eq:inner-estimate} in the numerator and
\eqref{eq:full-equivalent} in the denominator. For
$|j-z_n|\le A\sigma_n$, the ratio of the analytic prefactors tends to
$1$, while \eqref{eq:quadratic-window} gives the Gaussian exponent.
The remaining constant is
$\sqrt{\Delta_n}/\sqrt{2\pi B_n}=1/(\sqrt{2\pi}\sigma_n)$,
proving \eqref{eq:local-clt}.

The same estimates on windows of growing standardized width, together
with the curvature tail bound, imply conditional tightness and convergence
of interval probabilities. This proves the conditional central limit theorem.
If $M(p-1)>1$, then $\Prob(\One)\to1$ by the cutoff theorem. For every
bounded continuous $h$, conditioning changes
$\E h((J_n-z_n)/\sigma_n)$ by at most
$2\|h\|_\infty\Prob(\One^c)$, which tends to zero. Thus the
unconditional weak limit is the same. The variance scale in
\eqref{eq:variance} follows from \eqref{eq:BD-asymptotics}.
\end{proof}

\subsection{Why the leading action scale is not an adequate center}

The approximation $z_n\sim(p/(p-1))n/\log n$ determines only the first
scale. Even the next deterministic correction is generally much larger
than $\sigma_n\asymp\sqrt n/\log n$. For the bare equation
\eqref{eq:bare-r}, for example,
\[
 pr_n=(p-1)\log n-p\log\log n+O(1),
\]
so replacing $r_n$ by its first term shifts $j_n^0$ by order
$n\log\log n/(\log n)^2$. Dividing that shift by $\sigma_n$ tends to
infinity. A central limit theorem written with only the leading scale as
its center would therefore usually be false.

The coupled equations \eqref{eq:saddle} avoid this problem: they define a
center accurate to the Gaussian scale without requiring a manually expanded
series of logarithms. They are not an assertion that $z_n$ equals the
finite-$n$ mean exactly. The theorem is the centered distributional limit,
with the separately stated conditional local refinement.

\section{Explicit equivalents and the quadratic correction}
\label{sec:explicit}

\subsection{The bare saddle}

Consider first the comparison sum
\begin{equation}
 T_n=\sum_{j=1}^n\frac{(a j^p)^{n-j}}{(n-j)!}.
 \label{eq:bare-sum}
\end{equation}
It corresponds to the same coefficient/action calculation with $G(t)=t$
and $R(t)=1$. The saddle equations reduce to
\[
 k=a j^pt,\qquad-\log t=\frac{pk}{j},\qquad j+k=n.
\]
Writing $r=k/j$ gives exactly \eqref{eq:bare-r}. Its left side is
strictly increasing from $0$ to $\infty$, so the positive solution is
unique. At this saddle,
\[
 t=e^{-pr},\quad f=k(1+pr),\quad
 B=k,\quad\Delta=(1+pr)^2-p(p-1)r^2
                  =1+2pr+pr^2.
\]
The preceding saddle proof therefore gives $T_n\sim S_n$, with $S_n$
as in \eqref{eq:bare-S}.

This comparison has a simple interpretation: the slope at the giant is
kept, while every finite-core feedback correction is suppressed. Whether
that suppression is harmless depends on $p$.

\subsection{A controlled perturbation of the saddle}

\begin{proposition}[First core correction]
\label{prop:first-correction}
Under \eqref{eq:assumptions}, write
$\beta=u_2=\lambda_1+w_2$. If $p>3/2$, then
\begin{equation}
 u_n\sim S_n\exp\!\left(\frac{\beta(k_n^0)^2}{a(j_n^0)^p}\right).
 \label{eq:general-first-correction}
\end{equation}
In particular, for $p>2$ any finite perturbation satisfying
\eqref{eq:assumptions} has $u_n\sim S_n$.
\end{proposition}

\begin{proof}
Choose a fixed $M\ge\max(M_0,2)$ with $M(p-1)>1$. Then
$G_M(t)=t+\beta t^2+O(t^3)$ and $R_M(t)=1+O(t)$.
For a real $j$ in the coarse saddle interval put
$k=n-j$, $\Lambda=a j^p$, and $t_0=k/\Lambda$. The small solution of
$\Lambda tG'(t)=k$ has the uniform analytic expansion
\[
 t=t_0-2\beta t_0^2+O(t_0^3).
\]
Substitution into the exponent yields
\begin{align}
 f_G(j)
 &=k-k\log t_0+\beta k t_0+O(k t_0^2)\notag\\
 &=f_0(j)+\delta(j)+O(n t_0^2),
 \qquad\delta(j)=\frac{\beta(n-j)^2}{a j^p}.
 \label{eq:perturbed-exponent}
\end{align}
The analytic remainder can be differentiated with respect to $j$;
since $t_0'/t_0=-1/k-p/j=O(1/j)$, its first derivative is
$O(n t_0^2/j)$. Similarly $\delta'(j)=O(n t_0/j)$.

Let $j_0=j_n^0$ maximize $f_0$ and let $z$ maximize $f_G$.
Both lie in the same coarse interval, and its curvature has magnitude
comparable to $n/j_0^2$. The derivative estimates and the mean-value theorem
therefore give $z-j_0=O(j_0 t_0(j_0))$. Taylor's theorem then gives
\begin{equation}
 f_G(z)-f_0(j_0)=\delta(j_0)+O(n t_0(j_0)^2).
 \label{eq:peak-perturbation}
\end{equation}
Here the quadratic shift in the maximum, the change in $\delta$, and the
explicit remainder in \eqref{eq:perturbed-exponent} all have the displayed
order. The determinant prefactors have ratio $1+O(t_0(j_0))$, and
$R_M(\tau_n)=1+o(1)$. These estimates follow directly from
\eqref{eq:determinant}, because $z/j_0=1+O(t_0)$ and the core
and its first two derivatives differ from those of $t$ by $O(t)$ in
relative scale where appropriate.

Finally,
\[
 t_0(j_0)\asymp n^{1-p}(\log n)^p,\qquad
 n t_0(j_0)^2\asymp n^{3-2p}(\log n)^{2p}=o(1)
\]
when $p>3/2$. Combine \eqref{eq:peak-perturbation} with
Theorem~\ref{thm:saddle} and the cutoff theorem. This proves
\eqref{eq:general-first-correction}. For $p>2$, the retained correction
$\beta k_0t_0=O(n^{2-p}(\log n)^p)$ also tends to zero.
\end{proof}

\begin{corollary}[Quadratic feedback]
For the pure quadratic model,
\begin{equation}
 r_n(1+r_n)e^{2r_n}=an,
 \qquad
 u_n\sim
 \frac{\exp\!\left(\frac{nr_n}{1+r_n}(1+2r_n)
                      +(1+1/a)r_n^2\right)}
      {\sqrt{1+4r_n+2r_n^2}}.
 \label{eq:quadratic-full}
\end{equation}
The ratio $u_n/S_n$ tends to infinity faster than every fixed power of $n$.
\end{corollary}

\begin{proof}
Here $\beta=a+1$ and $k_0/j_0=r_n$, so the correction in
\eqref{eq:general-first-correction} is $(1+1/a)r_n^2$.
Also $r_n\sim\tfrac12\log n$. Thus
$\log(u_n/S_n)\sim(1+1/a)(\log n)^2/4$, which dominates every
fixed multiple of $\log n$.
\end{proof}

The order-$n$ logarithmic asymptotic cannot detect this factor because
$(\log n)^2=o(n)$. The quadratic example is therefore a concrete reason
to distinguish logarithmic control from multiplicative control.

\subsection{The hierarchy below the quadratic regime}

For $3/2<p<2$, the first correction in
\eqref{eq:general-first-correction} grows as a positive constant times
$n^{2-p}(\log n)^p$ in the pure model. The same two-core calculation still
gives a relative equivalent, because the next exponent scale tends to zero.
At $p=3/2$, that next scale is logarithmically divergent; a two-action core
is no longer enough. More generally, the small-action scales are
\begin{equation}
 n\tau_n^{d-1}\asymp
 n^{1-(d-1)(p-1)}(\log n)^{p(d-1)}.
 \label{eq:hierarchy}
\end{equation}
Theorem~\ref{thm:cutoff} handles the whole hierarchy by retaining an exact
finite analytic core instead of expanding each of these exponent scales by
hand. In the pure model, representative choices are
\begin{center}
\begin{tabular}{cc}
\toprule
Range of $p$ & Smallest sufficient $M$\\
\midrule
$p>2$ & $1$\\
$3/2<p\le2$ & $2$\\
$4/3<p\le3/2$ & $3$\\
$5/4<p\le4/3$ & $4$\\
\bottomrule
\end{tabular}
\end{center}
The core must grow as $p\downarrow1$. None of the fixed-$p$ arguments
above asserts a uniform theorem in that limit.

\section{Coefficient ratios and the least formal term}
\label{sec:least-term}

The multiplicative theorem also determines the scale at which the terms
of the divergent transseries stop decreasing. This is a statement about
formal terms, not yet an analytic remainder theorem.

\begin{proposition}[Ratio law throughout the superlinear range]
\label{prop:all-ratios}
Under \eqref{eq:assumptions}, for every fixed $p>1$,
\begin{equation}
 \frac{u_n}{u_{n-1}}\sim\frac1{\tau_n}
 \sim e^{pr_n}
 \sim a\left(\frac p{p-1}\right)^p
           \frac{n^{p-1}}{(\log n)^p},
 \label{eq:all-ratios}
\end{equation}
where $\tau_n$ comes from any sufficient core and $r_n$ solves
\eqref{eq:bare-r}.
\end{proposition}

\begin{proof}
Fix a sufficient core and extend its saddle smoothly to real $n$.
The implicit-function theorem applies because $\Delta_n>0$.
Differentiating the two stationary equations gives the useful exact identities
\begin{equation}
 \frac{\dd z_n}{\dd n}=\frac{C_n}{\Delta_n},\qquad
 \frac{\dd\log\tau_n}{\dd n}=-\frac{H_n}{\Delta_n}.
 \label{eq:saddle-motion}
\end{equation}
For completeness, at fixed $j$ the inner equation gives
$\partial_n\log t_j=1/B_j$, while
$\partial_j\log t_j=-C_j/B_j$ and
$\partial_n f_n'(j)=C_j/B_j$. Differentiating
$f_n'(z_n)=0$ proves the first identity; substituting it into the total
derivative of $\log t_{z_n}$ proves the second.

The scales already established imply
$z_n'=O(z_n/n)$ and $\tau_n'/\tau_n=O(1/n)$.
Differentiating the finite analytic expressions in
\eqref{eq:determinant} therefore gives
$\dd\log\Delta_n/\dd n=O(1/n)$. There is no dangerous cancellation
here, since $\Delta_n\asymp n^2/z_n^2$ and the differentiated terms
are individually $O(\Delta_n/n)$. Likewise
$\dd\log R(\tau_n)/\dd n=O(\tau_n/n)$.
The stationary exponent satisfies the exact envelope identity
\[
 \frac{\dd}{\dd n}
 \bigl(a z_n^pG(\tau_n)-(n-z_n)\log\tau_n\bigr)
 =-\log\tau_n.
\]
Thus $\dd\log\Acal_{n,M}/\dd n=-\log\tau_n+O(1/n)$.
Taking a difference over one unit and using $u_n\sim\Acal_{n,M}$ proves
the first equivalence in \eqref{eq:all-ratios}.
The full-core and bare radii have ratio tending to $1$: both are
asymptotic to the right side of \eqref{eq:saddlescales}. The last two
equivalences follow from \eqref{eq:bare-r}.
\end{proof}

\begin{corollary}[Least-term index]
\label{cor:least-term}
In the pure model, let $N(q)$ be any integer minimizing $u_nq^n$ over
$n\ge1$, for $q>0$. Such a minimum exists, and as $q\downarrow0$,
\begin{equation}
 \boxed{\quad
 N(q)\sim a^{-1/(p-1)}q^{-1/(p-1)}
       \left(\frac{\log(1/q)}p\right)^{p/(p-1)}.
 \quad}
 \label{eq:least-index}
\end{equation}
Equivalently, for $q=e^{-x}$,
$N(e^{-x})\sim a^{-1/(p-1)}(x/p)^{p/(p-1)}e^{x/(p-1)}$.
\end{corollary}

\begin{proof}
The coefficients are positive, and their $n$th roots tend to infinity by
Lemma~\ref{lem:macro}; hence $u_nq^n\to\infty$ for each fixed $q>0$,
so a minimum exists. Its index tends to infinity as $q\downarrow0$, because
for every fixed $n$ the ratio $q u_{n+1}/u_n$ then tends to zero.

Let $x=\log(1/q)$ and $r=x/p$. The real index at which the \emph{bare}
radius equals $q$ is exactly
\begin{equation}
 \nu(q)=a^{-1/(p-1)}r^{1/(p-1)}(1+r)q^{-1/(p-1)}.
 \label{eq:least-bare-index}
\end{equation}
The smooth comparison $e^{pr_n}$ in \eqref{eq:all-ratios} is increasing
and regularly varying with index $p-1$. For any fixed $\delta>0$, the
ratios $q u_n/u_{n-1}$ are therefore below $1$ on all sufficiently large
indices up to $(1-\delta)\nu(q)$, and above $1$ from
$(1+\delta)\nu(q)$ onward. This follows uniformly from the eventual
relative bounds in \eqref{eq:all-ratios} and monotonicity of the smooth
comparison; the finitely many small indices also have ratios below $1$
for small enough $q$. Every minimizing index lies between these two
bounds. Letting $\delta\downarrow0$ gives $N(q)\sim\nu(q)$, whose
leading form is \eqref{eq:least-index}.
\end{proof}

\begin{remark}[A boundary that must not be crossed without another proof]
The corollary locates the least coefficient term. It does not show that a
sectorial analytic solution differs from an optimally truncated series by
at most that term, or by any fixed multiple of it. Such a conclusion needs
uniform remainder estimates for the actual function. Keeping these two
statements separate is particularly important for the rapidly growing
feedback kernel.
\end{remark}

\section{Sharp reversion in the superquadratic regime}
\label{sec:inverse}

\subsection{A rapid-coefficient convolution lemma}

We isolate the signed part of the argument. Positivity was essential in
the forward cutoff proof, so that proof must not simply be reused after
formal reversion.

\begin{lemma}[Endpoint domination and reversion]
\label{lem:rapid-inverse}
Let $A(q)=q+\sum_{n\ge2}a_nq^n$, with $a_n>0$. Suppose there is a
smooth, eventually convex function $F$ such that
\[
 a_n\sim e^{F(n)},\qquad
 \frac{a_n}{a_{n-1}}\sim\mathcal R_n,
 \qquad
 \mathcal R_n\sim c\frac{n^\rho}{(\log n)^\kappa},
 \quad c>0,\quad\rho>1.
\]
For $A^{\langle-1\rangle}(q)=q+\sum_{n\ge2}\alpha_nq^n$,
\begin{equation}
 \alpha_n=-a_n+(n+1)a_2a_{n-1}+o(na_{n-1}).
 \label{eq:rapid-inverse}
\end{equation}
\end{lemma}

\begin{proof}
Set $B(q)=A(q)/q-1=\sum_{m\ge1}b_mq^m$, so $b_m=a_{m+1}$.
We first show
\begin{equation}
 [q^N]B(q)^2\sim2b_1b_{N-1},\qquad
 [q^N]B(q)^2\le Cb_{N-1}\quad(N\ge2)
 \label{eq:convolution}
\end{equation}
for a fixed constant $C$.
The summands with indices $1$ and $N-1$ give $2b_1b_{N-1}$.
For each fixed $j\ge2$,
\[
 \frac{b_jb_{N-j}}{b_{N-1}}=O(\mathcal R_N^{-(j-1)})=o(1).
\]
To bound the middle uniformly, choose a large fixed integer $L$ beyond
the convexity threshold of $F$. On $L\le j\le N-L$, the function
$F(j+1)+F(N-j+1)$ is convex, so its maximum occurs at the endpoints.
The comparability $a_n\asymp e^{F(n)}$ gives
\[
 \sum_{j=L}^{N-L}\frac{b_jb_{N-j}}{b_{N-1}}
 \le C_L N\frac{a_{N-L+1}}{a_N}
 =O(N\mathcal R_N^{-(L-1)})=o(1)
\]
when $L$ is sufficiently large. The finitely many endpoint terms between
$2$ and $L-1$ are already negligible. This proves the first assertion
of \eqref{eq:convolution}; increasing $C$ covers the finitely many small
values of $N$ and gives its second assertion.

Induction by convolution now gives, for $N\ge\ell\ge1$,
\begin{equation}
 [q^N]B(q)^\ell\le C^{\ell-1}b_{N-\ell+1}.
 \label{eq:iterated-convolution}
\end{equation}
Indeed, convolving the induction hypothesis with $B$ reduces the next
step to \eqref{eq:convolution} with total index $N-\ell+2$.

Lagrange inversion for $A^{\langle-1\rangle}$ yields
\begin{equation}
 \alpha_n=\frac1n\sum_{\ell=1}^{n-1}(-1)^\ell
 \binom{n+\ell-1}{\ell}[q^{n-1}]B(q)^\ell.
 \label{eq:inverse-binomial}
\end{equation}
The term $\ell=1$ is $-a_n$, while \eqref{eq:convolution} makes the term
$\ell=2$ equal to $(n+1)a_2a_{n-1}(1+o(1))$.
We bound all remaining terms absolutely, so no unproved cancellation
estimate is used.

For $3\le\ell\le n/2$, the ratio hypothesis, uniformly on indices
between $n/2$ and $n$, gives
\[
 \frac{a_{n-\ell+1}}{a_n}\le
 \left(\frac{C_1}{\mathcal R_n}\right)^{\ell-1}.
\]
Also $n^{-1}\binom{n+\ell-1}{\ell}\le
2(2n)^{\ell-1}/\ell!$. Thus the normalized absolute sum of these terms
is $O((n/\mathcal R_n)^2)$. Because $n/\mathcal R_n\to0$, this is
$o(n/\mathcal R_n)$, or $o(na_{n-1}/a_n)$.

For $\ell>n/2$, the numerator coefficient in
\eqref{eq:iterated-convolution} is at most a constant times
$a_{\lfloor n/2\rfloor+1}$. The ratio hypothesis implies
\[
 \log\frac{a_{\lfloor n/2\rfloor+1}}{a_n}
 =-\tfrac\rho2 n\log n+O(n\log\log n).
\]
This beats the $e^{O(n)}$ bound for the binomial factors and $C^\ell$,
and also the number of summands. Their contribution is smaller than
$o(na_{n-1})$. Combining the three ranges proves
\eqref{eq:rapid-inverse}.
\end{proof}

\subsection{Checking the hypotheses for the feedback coefficients}

\begin{proof}[Proof of the reversion and ratio parts of Theorem~\ref{thm:explicit}]
For $p>2$, Proposition~\ref{prop:first-correction} gives $u_n\sim S_n$.
Extend the positive solution $r_n$ of \eqref{eq:bare-r} smoothly to real
$x>0$. Implicit differentiation gives
\begin{equation}
 r'(x)=\frac{p-1}
 {x\left(p+1/r(x)+(p-1)/(1+r(x))\right)}
 \sim\frac{p-1}{px}.
 \label{eq:r-derivative}
\end{equation}
Let $F(x)=\log S_x$. The derivative of the optimized bare exponent with
respect to $x$ is $-\log t=pr(x)$; the derivative in the optimizing
action vanishes. The logarithmic determinant contributes only
$O(1/(x\log x))$. Consequently
\begin{equation}
 F'(x)=pr(x)+O\!\left(\frac1{x\log x}\right),\qquad
 F''(x)\sim\frac{p-1}{x}>0.
 \label{eq:F-derivatives}
\end{equation}
The second estimate follows either by differentiating the explicit
$r$-formula once more or directly from \eqref{eq:r-derivative}; its absolute error
is $O(1/(x\log x))$, hence its relative error is $O(1/\log x)$.
Thus $F$ is eventually convex.

The mean-value theorem and \eqref{eq:F-derivatives} show
$S_n/S_{n-1}\sim e^{pr_n}$. The same is true for $u_n/u_{n-1}$ because
$u_n/S_n\to1$. From \eqref{eq:bare-r},
\begin{equation}
 e^{pr_n}
 =\frac{a n^{p-1}}{r_n(1+r_n)^{p-1}}
 \sim a\left(\frac p{p-1}\right)^p
       \frac{n^{p-1}}{(\log n)^p}.
 \label{eq:rapid-ratio-scale}
\end{equation}
Lemma~\ref{lem:rapid-inverse} applies with $\rho=p-1>1$ and
$\kappa=p$, proving \eqref{eq:inverse}. In particular,
\begin{equation}
 -\frac{v_n}{u_n}
 =1-\frac{(n+1)u_2}{e^{pr_n}}
       +o\!\left(\frac n{e^{pr_n}}\right),
 \label{eq:inverse-relative}
\end{equation}
and $v_n<0$ for all sufficiently large $n$.
The forward parts of the theorem were proved in
Proposition~\ref{prop:first-correction} and its quadratic corollary.
\end{proof}

\subsection{The exact point where the proof stops}

The small parameter in the inverse proof is
\[
 \epsilon_n:=\frac n{\mathcal R_n}
 \asymp n^{2-p}(\log n)^p.
\]
It tends to zero precisely in the superquadratic range considered here.
At $p=2$ it instead grows logarithmically, and the higher powers in
\eqref{eq:inverse-binomial} cannot be discarded by this argument.
The forward cutoff theorem remains valid there, but its positivity is
lost on reversion. We do not extend \eqref{eq:inverse} into that range by
analogy or by extrapolating numerical ratios.

% ed. (2026-09-29): editorial cross-reference to the weighted-type package (batch 48) and to the two batch-49 claimed proofs
\begin{ednote}
The exponential ($n$th-root) scale of the inverse coefficients for
$1<p\le2$ is settled in limsup form by the later package \cite{ed:swt},
which takes two forward results of \cite{feedback} as inputs: for
$\lambda_j=aj^p$ its Theorem~\texttt{thm:refinedfeedback} gives
$\limsup_n\bigl(|v_n|\log(n+\ee)^{pn}/(n!)^{p-1}\bigr)^{1/n}
=a(p/(p-1))^p$, the forward value. For $p>2$ this value already follows
from \eqref{eq:inverse}. At $p=2$ the factors separating the conjectured
inverse equivalent from $u_n$ are $\ee^{o(n)}$, so that result neither
proves nor contradicts Conjecture~\ref{conj:quadratic-inverse}; for the
claimed proofs of the conjecture see the note after it.
\end{ednote}

\section{Algorithms, exact checks, and numerical diagnostics}
\label{sec:computations}

\subsection{Positive coefficient recurrence}

Let $E_{j,m}=[q^m]e^{\lambda_j\Uhat(q)}$ and $E_{j,0}=1$.
Differentiation of the exponential gives the triangular recurrence
\begin{equation}
 u_n=\sum_{j=1}^n w_jE_{j,n-j},\qquad
 E_{j,m}=\frac{\lambda_j}{m}
              \sum_{k=1}^m k u_k E_{j,m-k}\quad(m\ge1).
 \label{eq:numeric-recurrence}
\end{equation}
For nonnegative input, every term in these sums is nonnegative. Computing
in logarithms with a log-sum-exp operation therefore avoids both overflow
and subtractive cancellation. Only the triangle $j+m\le N$ is needed
through order $N$. The direct implementation takes $O(N^3)$ arithmetic
operations and $O(N^2)$ storage; this is not a bit-complexity bound for
large-integer arithmetic.

For integer slopes, factorial normalization avoids rational intermediate
values. Put $b_n=n!u_n$ and $\mathcal E_{j,m}=m!E_{j,m}$. Then
\begin{align}
 b_n&=\sum_{j=1}^n w_j\frac{n!}{(n-j)!}\mathcal E_{j,n-j},\\
 \mathcal E_{j,m}
 &=\lambda_j\sum_{k=1}^m\binom{m-1}{k-1}
                    b_k\mathcal E_{j,m-k}.
 \label{eq:integer-recurrence}
\end{align}
The supplied exact implementation uses $w_j=1$, integer $a$ and integer
$p$. It computes the inverse in the same factorial normalization by the
partial Bell-polynomial recurrence for the identity
$\Uhat(\Vhat(q))=q$.

\subsection{Evaluating the finite-core formula}

For numerical evaluation one need not first compute a long Taylor expansion
of $G_M$. Solve simultaneously for $(t,u,z)$:
\begin{equation}
 u=\Phi_M(t,u),\qquad
 \log t+apz^{p-1}u=0,\qquad
 z+a z^ptG_M'(t)=n.
 \label{eq:numeric-saddle}
\end{equation}
The required derivatives are finite expressions:
\begin{align}
 G_M'&=\frac{\Phi_t}{1-\Phi_u},\\
 G_M''&=\frac{\Phi_{tt}+2\Phi_{tu}G_M'
                              +\Phi_{uu}(G_M')^2}{1-\Phi_u}.
 \label{eq:core-derivatives}
\end{align}
All partial derivatives are evaluated at $(t,u)$. The pure-model program
uses logarithms of the three positive unknowns, initializes with the bare
saddle, and checks the regular branch $\Phi_u<1$, a positive determinant,
and the numerical residual. These are floating-point checks, not interval
certificates. The existence theorem ensures the intended small branch for
sufficiently large $n$; it does not supply an explicit universal starting
order for a generic nonlinear root finder.

Each evaluation of the core and its derivatives uses only $M$ summands.
Thus the finite-core formula is useful even when storing the exact
coefficient itself is impractical. A larger sufficient core can improve
finite-order accuracy substantially without changing the limit theorem.

\subsection{What was actually checked}

The package records exact coefficients and formal inverse coefficients
through degree $300$ for $(a,p)=(1,2)$ and $(1,3)$. It contains:
\begin{itemize}
\item 36 independent comparisons of the positive recurrence with the
partition formula: degrees $1$ through $18$ in the two models;
\item 96 independent one-tail marking comparisons: degrees $1$ through
$16$, cutoffs $M=1,2,3$, in both models;
\item 598 exact inverse-composition residual checks: degrees $2$ through
$300$ in the two models.
\end{itemize}
All of these checks passed. The inverse coefficients are obtained from a
recurrence enforcing the composition identity; the residual checks are not
presented as an independent second proof of formal inversion. By contrast,
the partition enumeration and the one-tail enumeration use different finite
algorithms from the differential recurrences.

Separately, the log-domain positive recurrence was run through degree
$1000$ in both models. Comparing it with the exact coefficients through
$300$ gave maximum absolute discrepancies in $\log u_n$ of approximately
$6.83\times10^{-13}$ for $p=2$ and $9.10\times10^{-13}$ for $p=3$.
These agreement figures are diagnostics in ordinary floating-point
arithmetic, not rigorous error bars.

\subsection{The effect of the finite-core cutoff}

Table~\ref{tab:cutoff} compares the computed coefficient $u_{1000}$ with
$\Acal_{1000,M}$ from \eqref{eq:full-equivalent}, for $a=1$.
The quadratic $M=1$ core is insufficient and has a large missing factor.
The smallest sufficient core is not necessarily the most accurate choice
at a moderate degree; adding a few more small actions improves the numerical
approximation markedly.

\begin{table}[htbp]
\centering
\caption{Finite-core diagnostics at $n=1000$, $a=1$. The ratio uses the
log-domain recurrence; no interval certification is asserted.}
\label{tab:cutoff}
\begin{tabular}{rrrr}
\toprule
$p$ & $M$ & $u_n/\Acal_{n,M}$ & saddle action $z_n$\\
\midrule
2 & 1 & 373.908262191 & 292.05833 \\
2 & 2 & 1.051182297 & 290.22703 \\
2 & 3 & 1.000393348 & 290.19821 \\
2 & 4 & 0.999965103 & 290.19784 \\
3 & 1 & 1.045788395 & 235.38637 \\
3 & 2 & 0.999941539 & 235.37627 \\
3 & 3 & 0.999938904 & 235.37627 \\
3 & 4 & 0.999938904 & 235.37627 \\
\bottomrule
\end{tabular}
\end{table}

For the pure quadratic model at $n=1000$, the ratio to the bare formula is
approximately $124136.898$, while the ratio to the explicit quadratic
correction \eqref{eq:quadratic} is $1.19599044$. The exact four-action
core improves the ratio to approximately $0.99996510$.
The large bare error is not a finite-precision accident: its divergence is
proved by the quadratic corollary. Conversely, the small observed
four-core error is not a proven finite-$n$ guarantee.

\subsection{Inverse diagnostics and a deliberately separate conjecture}

The middle column of Table~\ref{tab:inverse} shows $-v_n/u_n$ from exact
integer-normalized coefficients. The superquadratic ratios move toward the
proved limit $1$, albeit slowly. The quadratic ratios are far smaller.
The last column compares with the conjectural extension discussed in
Section~\ref{sec:research}; that column is exploratory, not a theorem check.

\begin{table}[htbp]
\centering
\caption{Exact-coefficient inverse diagnostics, $a=1$. Here
$\widetilde V_n=-S_n\exp(-2k_n^0/(j_n^0)^{p-1})$ is an exploratory
comparison; its quadratic asymptotic status is conjectural.}
\label{tab:inverse}
\begin{tabular}{rrrr}
\toprule
$p$ & $n$ & $-v_n/u_n$ & $v_n/\widetilde V_n$\\
\midrule
2 & 50 & 0.000519898 & 0.649597 \\
2 & 100 & 0.000121097 & 0.757155 \\
2 & 200 & 0.000018894 & 0.835910 \\
2 & 300 & 0.000005371 & 0.870969 \\
3 & 50 & 0.571305779 & 0.978478 \\
3 & 100 & 0.667665290 & 0.992327 \\
3 & 200 & 0.751191069 & 0.997274 \\
3 & 300 & 0.793480085 & 0.998512 \\
\bottomrule
\end{tabular}
\end{table}

No finite collection of these comparisons proves the cutoff, the saddle
limit, the central limit theorem, or the inverse asymptotic. Their role is
to detect normalization errors, missing factors, and implausible finite-order
behavior in formulas proved by other means.

\section{Further research questions and a concrete inverse conjecture}
\label{sec:research}

\subsection{Quadratic and subquadratic inverse asymptotics}

The most immediate continuation is a finite-core theorem for the signed
inverse coefficients when $1<p\le2$. There is an exact identity suggesting
a route, but the needed cancellation control has not been proved here.

\begin{proposition}[Exact inverse response to one marked tail]
\label{prop:inverse-response}
Let $Q=G_M^{\langle-1\rangle}$. Fix $j>M$ and solve
\[
 U_s(q)=\Phi_M(q,U_s(q))+s q^j e^{a j^pU_s(q)},
 \qquad V_s=U_s^{\langle-1\rangle}.
\]
Then the formal derivative of the inverse at $s=0$ is
\begin{equation}
 \left.\partial_sV_s(u)\right|_{s=0}
 =-Q'(u)R_M(Q(u))Q(u)^j e^{a j^p u}.
 \label{eq:inverse-response}
\end{equation}
\end{proposition}

\begin{proof}
The direct response at $s=0$ is
$T_j(q)=R_M(q)q^je^{a j^pG_M(q)}$ by Lemma~\ref{lem:marking}.
Differentiate $U_s(V_s(u))=u$ at $s=0$:
$T_j(Q(u))+G_M'(Q(u))\partial_sV_s(u)|_{s=0}=0$.
Use $Q'=1/G_M'(Q)$ and $G_M(Q(u))=u$.
\end{proof}

For $M\ge2$ in the pure model,
$Q(u)=u-(a+1)u^2+O(u^3)$. Thus the large power $Q(u)^j$ in
\eqref{eq:inverse-response} suggests an exponential penalty of leading
form $\exp(-(a+1)ju)$, not the positive forward factor
$\exp((a+1)ku)$. At the bare saddle this suggests the following candidate.

\begin{conjecture}[Quadratic inverse equivalent]
\label{conj:quadratic-inverse}
For the pure model with $p=2$ and every fixed $a>0$,
\begin{equation}
 v_n\sim-S_n\exp\!\bigl(-(1+1/a)r_n\bigr).
 \label{eq:inverse-conjecture}
\end{equation}
\end{conjecture}

The exact inverse-response identity does not prove that the one-tail
response exhausts the inverse coefficient. That is the missing step.
The quadratic data in Table~\ref{tab:inverse} are consistent with the
candidate, but do not establish its constant, sign for all large $n$, or
its limit. For $p=3$, convergence of the table's last-column ratio to $1$
already follows from $v_n\sim-S_n$, because that column's extra exponential
factor tends to $1$; no refined rate follows from it.
A useful attack would reorganize the signed Lagrange formula into an
analytic core before bounding its remaining tail, rather than estimate
individual inverse-composition powers absolutely as in the $p>2$ proof.

% ed. (2026-09-29): editorial note recording the two independent claimed proofs of conj:quadratic-inverse (batch 49); the conjecture itself is unchanged
\begin{ednote}
Conjecture~\ref{conj:quadratic-inverse} remains stated here as this
article's conjecture. Two later packages, filed in batch~49 (see
\path{docs/incoming/README.md}), each claim a proof of it; neither proof
has been independently reviewed, and neither package cites the other.
They are \cite{ed:qef}, Theorem~\texttt{thm:main}, and \cite{ed:sqf},
Theorem~\texttt{thm:main}. Both use $r_n$ and $S_n$ of
\eqref{eq:bare-r} and \eqref{eq:bare-S} and follow the route suggested
above, resumming an analytic inverse core before bounding the signed tail:
the first merges multiple large actions in a positive majorant, the second
uses an exact quadratic merger defect. Their exact inverse coefficients
agree at every common degree (through $220$ for $a=1$ and $180$ for
$a=2$). Neither hypothesis set contains the other. The first allows
finitely many nonnegative changes of weights and slopes and eventual
slopes $aj^2+dj+e$, with factor $\exp((d-\mu)r_n/a)$,
$\mu=\lambda_1+w_2$; the second allows arbitrary real changes of finitely
many weights and slopes but an exact tail $aj^2$, with factor
$\exp(-(\lambda_1+w_2)r_n/a)$. In the pure model both factors reduce to
$\exp(-(1+1/a)r_n)$, which is \eqref{eq:inverse-conjecture}.
% ed. (2026-09-29): added on filing batch 50; a third independent claimed proof.
A third, \path{Signed_Condensation_Sharp_Quadratic_Reversion/} (batch~50),
claims it under the first package's hypotheses (its
Theorem~\texttt{thm:main}) and adds the least formal inverse term; its
exact coefficients agree with both (through degree $240$ for $a=1$ and
$160$ for $a=2$). It has not been independently reviewed.
\end{ednote}

\subsection{Regularly varying feedback and critical logarithmic boundaries}

Replace the pure tail by
$\lambda_j\sim a j^p(\log j)^\gamma$, with controlled derivatives or
finite differences. The leading radius is then expected to have scale
$n^{1-p}(\log n)^{p-\gamma}$. The corresponding omitted-action diagnostic
would be
\[
 n\tau_n^M\asymp
 n^{1-M(p-1)}(\log n)^{M(p-\gamma)}.
\]
At $M(p-1)=1$, this predicts three different possibilities: a vanishing
correction for $\gamma>p$, a finite nonzero correction for $\gamma=p$,
and a divergent correction for $\gamma<p$. Establishing a theorem at that
boundary requires a sharper replacement for the power-tail transfer bound,
and potentially a Poisson correction when the scale tends to a constant.
The power-tail theorem does not automatically cover these alternatives.

% ed. (2026-09-29): editorial cross-reference to the Poisson-layer package (batch 48), which answers part of this question
\begin{ednote}
The later package \cite{ed:plb} obtains such a Poisson correction at
$M(p-1)=1$ by letting the exponent approach the boundary,
$p_n=1+1/M+((M+1)\log\log n+\theta)/(M\log n)$
(Theorem~\texttt{thm:window}); the $(\log j)^\gamma$ tails proposed here
are not treated there and remain open.
\end{ednote}

\subsection{Joint laws for the finite background cloud}

The present local theorem concerns the large action after conditioning on
one tail. Determine the joint law of the numbers of actions
$2,\ldots,M$ and their dependence on $J_n$. Natural scales are
$n\tau_n^{d-1}$, but their exact means and covariances depend on the core's
implicit response. One can introduce finitely many marking variables in
$\Phi_M$ and differentiate the saddle equations. The desired theorem
would give Gaussian limits for diverging counts, Poisson limits for bounded
counts, and a quantitative total-variation approximation where appropriate.
It should explicitly distinguish this coefficient model from the fixed
product-weight random-allocation models in \cite{jjs,janson}.

% ed. (2026-09-29): editorial cross-reference to the two batch-48 packages that answer this question in special cases
\begin{ednote}
Two later packages answer this question in special cases. The package
\cite{ed:plb} treats the single count $N_{M+1,n}$ of first omitted actions
near the cutoff boundary: Poisson limits in total variation for bounded
intensity, Gaussian limits for diverging intensity, and independence from
the largest action (Theorems~\texttt{thm:probability} and
\texttt{thm:joint}). The package \cite{ed:mic} treats $p=2$, $a=1$: the
number of actions equal to two is close in total variation to a Poisson law
of mean $r_n^2$, jointly with the largest action
(Theorem~\texttt{thm:probability}), and with high probability every action
other than the largest is one or two. For other $p$, and for $p=2$ with $a\ne1$, the joint
law of the counts of actions $2,\ldots,M$ remains open.
\end{ednote}

\subsection{Uniformity as the growth exponent approaches one}

The minimum sufficient core diverges as $p\downarrow1$. Moreover, when
$(p-1)\log n$ is bounded, the proposed giant size is no longer small
relative to $n$. Both the fixed-core analytic estimates and the coarse
localization need replacement. Develop a double-scaling theorem in which
$p=p_n\downarrow1$ and determine the transition between a macroscopic
fraction of total action and the $n/\log n$ scale. A uniform result here
would connect the convergent linear-feedback boundary with the divergent
superlinear regime, rather than merely taking a limit in a nonuniform
formula.

\subsection{Quantified errors and certified saddle computation}

The current asymptotic theorems have $o(1)$ relative errors without explicit
universal onset values. Convert the mass-transfer inequality, the Cauchy
arc bounds, and the action-saddle Taylor estimates into explicit constants.
Then combine interval evaluation of the finite sums in
\eqref{eq:numeric-saddle} with a bracket for the unique coarse saddle.
The target is a computable pair of bounds
$L_n\le u_n/\Acal_{n,M}\le U_n$ with $L_n,U_n\to1$ and a
specified domain of validity. Floating-point residuals alone, even when
small, do not provide that certificate.

\subsection{Higher-order expansions and accuracy-dependent cores}

A complete expansion should combine higher Cauchy-saddle terms, a discrete
Edgeworth expansion for the action sum, and omitted-background corrections.
For a target algebraic accuracy, the natural first estimate is to choose a
core with $n\tau_n^M$ below that accuracy scale. This suggests stricter
cutoffs than $M(p-1)>1$, but the logarithmic factors and the analytic-saddle
errors must be tracked together. Determine the smallest core for a prescribed
relative error $O(n^{-K})$, and whether the resulting expansion can be
organized in a reusable symbolic coefficient calculus.

\subsection{From least formal terms to actual analytic remainders}

Corollary~\ref{cor:least-term} determines a formal truncation scale for every
$p>1$. The earlier feedback article constructs left-sector analytic
solutions \cite{feedback}, but a useful continuation would establish
uniform remainder estimates strong enough to compare those solutions with
truncation at that scale. Directional summability or acceleration requires
further analytic growth information; neither follows from a coefficient
ratio law alone. A precise target is a sectorial remainder bound with an
explicit action-dependent exponent, followed by an analysis of the Stokes
ambiguity between different realizations when such ambiguity exists.

\subsection{A staged formalization in Lean}

The shortest useful formalization path begins with finite algebra rather
than the full saddle theorem. Formalize the positive Lagrange coefficient
identity, its finite multiplicity form, and the marked one-tail response;
then the convex mass-transfer inequality and the preimage-count bound.
These are separate from the analytic implicit-function and contour estimates.
The existing transseries support and asymptotic-scale library provides
context, but no claim is made that it already supplies this entire analytic
chain. A later stage can package shrinking-radius saddle estimates with
uniform hypotheses, followed by the lattice Gaussian summation and signed
inverse lemma. This decomposition would make a partially completed
formalization useful without labeling the unformalized remainder as checked.

\section{Conclusion}

For power feedback, the passage from logarithmic estimates to relative
asymptotics is controlled by a finite but sharply determined amount of
background data. The correct reduction is not simply to one large action,
but to one large action dressed by a sufficient analytic core. The exact
threshold is $M(p-1)>1$, and its strict boundary is forced by an explicit
positive two-tail contribution.

Once this reduction is established, two ordinary saddle calculations give
the full coefficient factor and a moving Gaussian law. They also expose the
superpolynomial factor hidden in the quadratic logarithmic estimate and
locate the least formal term. A separate rapid-coefficient argument handles
formal reversion when $p>2$. The quadratic inverse, quantitative analytic
remainders, and uniform transition at $p=1$ remain focused next problems,
not conclusions silently assumed by the present results.

\appendix
\section{Proof dependencies and scope audit}
\label{app:audit}

The proof order avoids circular use of the cutoff theorem. First the positive
coefficient formula and coarse localization are proved. The upper side of
the cutoff uses a combinatorial transfer. Independently, the saddle formula
is proved for $w_{n,M}$ for every fixed core, including insufficient ones.
That independent formula is then used to prove failure below the cutoff.
Only after both sides are available is the unconditional fluctuation
statement obtained.

\begin{center}\small
\begin{tabular}{>{\raggedright\arraybackslash}p{0.25\textwidth}>{\raggedright\arraybackslash}p{0.40\textwidth}>{\raggedright\arraybackslash}p{0.25\textwidth}}
\toprule
Result & Main ingredients & Proven range\\
\midrule
Positive coefficients and one-tail identity & Formal Lagrange inversion;
analytic implicit core; marking & Assumptions \eqref{eq:assumptions}\\[4pt]
Macroscopic localization & Convexity, Stirling bounds, one-dimensional
strict maximum & Every fixed $p>1$\\[4pt]
Sufficient core & Slope transfer and explicit preimage count &
$M(p-1)>1$\\[4pt]
One-tail coefficient equivalent & Uniform Cauchy saddle; moving lattice
saddle & Every fixed $M\ge M_0$\\[4pt]
Necessary core & Exact two-mark response and saddle multiplier &
Failure for $M(p-1)\le1$\\[4pt]
Gaussian laws & Saddle window and tail control & Local conditional;
weak unconditional for sufficient $M$\\[4pt]
First explicit correction & Analytic perturbation of optimized exponent &
$p>3/2$\\[4pt]
Ratio and least-term index & Smooth saddle motion; regular variation &
$p>1$; least-term statement for pure model\\[4pt]
Sharp signed reversion & Endpoint convolution; absolute high-power bound &
Pure model, $p>2$\\[4pt]
Quadratic inverse equivalent & Exact response plus unproved tail dominance &
Conjecture only\\
\bottomrule
\end{tabular}
\end{center}

The ordinary analytic implicit-function theorem, Cauchy's coefficient formula,
Stirling's estimate, elementary convexity, and the Gaussian integral are
classical tools used in the proofs. Numerical agreement is not an additional
axiom or a replacement for any of them. No theorem in this document has
been compiled as Lean code. No repository file or branch was modified in
preparing the package.

The main claims are stated for nonnegative weights and slopes with an
\emph{exact power tail after finitely many exceptions}. Signed primitive
weights, an arbitrary regularly varying tail, a growing core, or a
parameter $p$ depending on $n$ require additional arguments. In particular,
a theorem about one of those larger classes is not implicit in our use of
the word universality.

\section{Notation and normalization}
\label{app:notation}

\begin{center}\small
\begin{tabular}{>{\raggedright\arraybackslash}p{0.24\textwidth}>{\raggedright\arraybackslash}p{0.69\textwidth}}
\toprule
Symbol & Meaning\\
\midrule
$q$ & Formal expansion variable; $q=e^{-x}$ gives the exponential
transseries interpretation.\\[3pt]
$w_j,\lambda_j$ & Primitive weight and feedback slope of action $j$.\\[3pt]
$u_n,v_n$ & Ordinary coefficients of $\Uhat$ and its formal inverse.
They are not exponential-generating-function coefficients.\\[3pt]
$G_M,R_M$ & Analytic finite core and its linear response near zero.\\[3pt]
$W_M,w_{n,M}$ & Formal series and coefficient for exactly one action above
$M$.\\[3pt]
$K_n,E_n,J_n$ & Random number of parts, one plus excess action, and largest
part in the normalized positive Lagrange measure.\\[3pt]
$z_n,\tau_n$ & Real moving action saddle and its small coefficient radius.\\[3pt]
$k_n=n-z_n$ & Deterministic complement of the saddle action; not the random
number of parts $K_n$.\\[3pt]
$B_n,C_n,H_n,\Delta_n$ & Derivative quantities in
\eqref{eq:determinant}; $\sigma_n^2=B_n/\Delta_n$.\\[3pt]
$r_n,j_n^0,k_n^0,S_n$ & Bare saddle ratio, action, complement, and coefficient
approximation.\\[3pt]
$b_n=n!u_n$ & Integer normalization used only in the exact-computation
section and output files.\\[3pt]
$\sim$ & Ratio tends to one, except when explicitly applied to logarithms.
All displayed asymptotics fix the model parameters first.\\
\bottomrule
\end{tabular}
\end{center}

For the pure model the least sufficient core is $M_*(p)$ from
\eqref{eq:mincore}. With arbitrary initial perturbations the theorem is
applied with any $M\ge M_0$ also satisfying the strict cutoff. This ensures
that every action subjected to the tail transfer has precisely the assumed
power slope and unit weight.

\section{Reproduction and source provenance}
\label{app:reproduction}

The main TeX file is self-contained: its bibliography and numerical tables
are embedded, with no external figures or bibliography database required.
Compile it twice with \texttt{pdflatex}. The companion files are
\texttt{README.md}, \texttt{SOURCES.md}, \texttt{PROOF\_STATUS.md}, a
\texttt{Makefile}, three Python programs, and their recorded data.
A Python environment with NumPy and SciPy is needed only for numerical
saddle diagnostics and table generation. The exact program uses the standard
library alone.

\begin{verbatim}
python code/verify_exact.py --order 300
python code/diagnostics.py --order 1000 --powers 2 3
python code/report_tables.py
pdflatex -interaction=nonstopmode -halt-on-error article.tex
pdflatex -interaction=nonstopmode -halt-on-error article.tex
\end{verbatim}

The stored CSV files with integer-normalized coefficients are exact.
The stored logarithms, nonlinear-solver residuals, and ratios to saddle
formulas are floating-point data. Re-running the scripts may change last
printed digits or timings across platforms. The article's asymptotic
proofs do not depend on those digits or on any external computation service.

The repository motivation is tied to the pinned source identified in
Section~\ref{sec:context}. The literature check was targeted, not an
exhaustive priority search. In particular, the paper treats the existence of
related condensation and rapid-series methods as prior work, while proposing
the exact finite-core reduction and the associated model-specific
asymptotics as the contribution to be independently reviewed.

% ed. (2026-09-29): widest label 9 -> 99, since the editorial entries bring the list to thirteen
\begin{thebibliography}{99}
\addcontentsline{toc}{section}{References}

\bibitem{repo}
Vladimir Reshetnikov, \emph{ProveIt: Analysis/Transseries}, repository and
research documentation. Snapshot consulted: commit
\path{5804c7aff9955e6cbc09be2a71a295a07b238bc9},
29 September 2026.
\href{https://github.com/VladimirReshetnikov/ProveIt/tree/5804c7aff9955e6cbc09be2a71a295a07b238bc9/Analysis/Transseries}{Pinned transseries directory}.

\bibitem{feedback}
\emph{A sharp regularity classification for countable exponential-feedback
transseries}, research draft prepared for Vladimir Reshetnikov,
29 September 2026. ProveIt,
\path{Analysis/Transseries/docs/series-and-transseries/Exponential_Feedback_Regularity_Classification/article.tex}.
Source blob \path{c6a83d71973a12fe7ed970d2abc56e2b49e6808f}.
\href{https://github.com/VladimirReshetnikov/ProveIt/blob/5804c7aff9955e6cbc09be2a71a295a07b238bc9/Analysis/Transseries/docs/series-and-transseries/Exponential_Feedback_Regularity_Classification/article.tex}{Pinned TeX source}.

\bibitem{gessel}
Ira M. Gessel, \emph{Lagrange Inversion}, Journal of Combinatorial Theory,
Series A \textbf{144} (2016), 212--249.
\href{https://arxiv.org/abs/1609.05988}{arXiv:1609.05988}.
\href{https://doi.org/10.1016/j.jcta.2016.06.018}{doi:10.1016/j.jcta.2016.06.018}.

\bibitem{jkt}
Sabine Jansen, Tobias Kuna, and Dimitrios Tsagkarogiannis,
\emph{Lagrange inversion and combinatorial species with uncountable color
palette}, Annales Henri Poincar\'e (2021).
\href{https://arxiv.org/abs/2008.10862}{arXiv:2008.10862}.
\href{https://doi.org/10.1007/s00023-020-01013-0}{doi:10.1007/s00023-020-01013-0}.

\bibitem{jjs}
Svante Janson, Thordur Jonsson, and Sigurdur Orn Stef\'ansson,
\emph{Random trees with superexponential branching weights},
Journal of Physics A: Mathematical and Theoretical \textbf{44} (2011),
485002.
\href{https://arxiv.org/abs/1104.2810}{arXiv:1104.2810}.
\href{https://doi.org/10.1088/1751-8113/44/48/485002}{doi:10.1088/1751-8113/44/48/485002}.

\bibitem{janson}
Svante Janson, \emph{Simply generated trees, conditioned Galton--Watson
trees, random allocations and condensation}, Probability Surveys
\textbf{9} (2012), 103--252.
\href{https://arxiv.org/abs/1112.0510}{arXiv:1112.0510}.

\bibitem{bender}
Edward A. Bender, \emph{An Asymptotic Expansion for the Coefficients of
Some Formal Power Series}, Journal of the London Mathematical Society
(2) \textbf{9} (1975), 451--458.
\href{https://doi.org/10.1112/jlms/s2-9.3.451}{doi:10.1112/jlms/s2-9.3.451}.

\bibitem{edgar}
G. A. Edgar, \emph{Transseries for beginners}, Real Analysis Exchange
\textbf{35} (2010), 253--310.
\href{https://arxiv.org/abs/0801.4877}{arXiv:0801.4877}.

% ed. (2026-09-29): editorial bibliography entries for the later sibling packages cited in the editorial notes
\bibitem{ed:mic}
[Editorial addition] \emph{Microscopic Condensation in Exponential-Feedback
Transseries}, research article prepared for Vladimir Reshetnikov,
ProveIt, filed in batch~48,
\path{Analysis/Transseries/docs/series-and-transseries/Microscopic_Condensation_Exponential_Feedback/article.tex}.

\bibitem{ed:plb}
[Editorial addition] \emph{Poisson Layers at the Finite-Core Boundary of
Exponential-Feedback Transseries},
research article prepared for Vladimir Reshetnikov, ProveIt, filed in
batch~48,
\path{Analysis/Transseries/docs/series-and-transseries/Poisson_Layers_Finite_Core_Boundary/poisson_feedback_layers.tex}.

\bibitem{ed:swt}
[Editorial addition] \emph{Sharp Weighted Type Under Formal Reversion},
research article prepared for Vladimir Reshetnikov, ProveIt, filed in
batch~48,
\path{Analysis/Transseries/docs/series-and-transseries/Sharp_Weighted_Type_Formal_Reversion/article.tex}.

\bibitem{ed:qef}
[Editorial addition] \emph{Quadratic Exponential Feedback after Reversion:
Cancellation, Finite-Core Universality, and Sharp Borel Growth}, research
article prepared for Vladimir Reshetnikov, ProveIt, filed in batch~49,
\path{Analysis/Transseries/docs/series-and-transseries/Quadratic_Exponential_Feedback_After_Reversion/article.tex}.
Claimed proof, not independently reviewed.

\bibitem{ed:sqf}
[Editorial addition] \emph{Signed Quadratic-Feedback Inversion: A
Finite-Core Resolution and Sharp Entire Borel Growth}, research draft
prepared for Vladimir Reshetnikov, ProveIt, filed in batch~49,
\path{Analysis/Transseries/docs/series-and-transseries/Signed_Quadratic_Feedback_Inversion/article.tex}.
Claimed proof, not independently reviewed.

\end{thebibliography}
\end{document}
